\documentclass[aps,prd,twocolumn,superscriptaddress,nofootinbib]{revtex4-2}

\usepackage{amsmath}
\usepackage{mathtools}
\usepackage{amssymb}
\usepackage{amsthm}
\usepackage{booktabs}
\usepackage{tabularx}
\usepackage{array}
\usepackage{graphicx}
\usepackage{xcolor}
\usepackage{mathrsfs}

\newtheorem{theorem}{Theorem}[section]
\newtheorem{proposition}[theorem]{Proposition}
\newtheorem{lemma}[theorem]{Lemma}
\newtheorem{corollary}[theorem]{Corollary}
\newtheorem{conjecture}[theorem]{Conjecture}
\newtheorem{axiom}[theorem]{Axiom}
\theoremstyle{definition}
\newtheorem{definition}[theorem]{Definition}
\theoremstyle{remark}
\newtheorem{remark}[theorem]{Remark}
\newtheorem{claim}[theorem]{Claim}

\DeclareMathOperator{\poly}{poly}
\DeclareMathOperator{\SIZE}{SIZE}
\DeclareMathOperator{\SAT}{SAT}
\newcommand{\NP}{\mathrm{NP}}
\newcommand{\PP}{\mathrm{P}}
\newcommand{\BPP}{\mathrm{BPP}}
\newcommand{\coNP}{\mathrm{coNP}}

\usepackage[hyperfootnotes=false]{hyperref}
\hypersetup{
    pdftitle={An Entropic Perspective on P vs NP: Reformulation via Kolmogorov Complexity},
    pdfauthor={Lukas Geiger},
    pdfsubject={Entropic reformulation of P vs NP via resource-bounded Kolmogorov complexity},
    pdfkeywords={P vs NP, Kolmogorov complexity, algorithmic entropy, witness compression, computational complexity},
    colorlinks=true,
    linkcolor=black,
    urlcolor=blue,
    citecolor=black
}
\renewcommand*{\HyperDestNameFilter}[1]{en.#1}

\begin{document}

% ===================================================================
% TITLE PAGE
% ===================================================================

\title{An Entropic Perspective on P vs NP: Reformulation via Kolmogorov Complexity}

\author{Lukas Geiger}
\email{https://orcid.org/0009-0005-7296-1534}
\affiliation{Independent Researcher, Bernau, Germany}

\date{May 5, 2026}

\begin{abstract}
We present a unified framework for the well-known equivalence between the $\PP$ vs $\NP$ problem and resource-bounded Kolmogorov complexity, framing it as an \textit{entropic reformulation}: $\PP \neq \NP$ is equivalent to the absence of a polynomial-time witness selector whose outputs have only logarithmic resource-bounded conditional complexity. This equivalence has been known in various forms since the work of Orponen et al.\ (1994) and Fortnow (2000). This paper does not prove $\PP \neq \NP$; rather, it provides a unified information-theoretic presentation that identifies the precise structure of what must be shown and why classical barrier techniques are not obstructions to this approach. We introduce the concept of \textit{algorithmic positivity} -- an NP language is algorithmically positive if it admits a polynomially scaling witness-producing function whose output has low resource-bounded conditional Kolmogorov complexity $K^t(W(x) \mid x) \leq O(\log|x|)$ -- and characterize the \textit{Entropic Separation Conjecture} (no NP-complete language is algorithmically positive) as \textit{equivalent} to $\PP \neq \NP$, establishing ESC as a reformulation of the P vs NP problem in information-theoretic language rather than an independent result. We conduct a systematic barrier analysis, arguing that arguments based on unbounded Kolmogorov complexity $K$ are not naturally subject to the three classical barriers (relativization, natural proofs, algebrization), while noting that the barrier status for the resource-bounded variant $K^t$ used in our core results remains to be fully established. The central characterization is the \textit{witness entropy gap}: $\PP = \NP$ would imply the existence of a canonical polynomial-time witness selector with $K^t(W(x) \mid x) = O(\log|x|)$, while $\PP \neq \NP$ implies that, for every fixed polynomial time bound and logarithmic constant, infinitely many instances admit no valid witness with such a short resource-bounded conditional description. The stronger near-maximal bound $K^t(w \mid x) \geq |w| - O(\log n)$ is treated below as an additional bridge hypothesis, not as a consequence of $\PP \neq \NP$. We identify the precise remaining challenge: converting this characterization into a proof requires a \textit{uniformity bridge} -- a new connection between algorithmic information theory and uniform complexity theory.
\end{abstract}

\keywords{P vs NP, Kolmogorov complexity, algorithmic entropy, reformulation, barrier analysis, witness compression, computational complexity}

\thanks{AI tools (Claude by Anthropic) were used for mathematical formalization, literature verification, and editorial refinement. All mathematical content was developed, verified, and approved by the author.}

\maketitle


% ===================================================================
% 1. INTRODUCTION
% ===================================================================
\section{Introduction}
\label{sec:intro}

\subsection{The problem}

The $\PP$ vs $\NP$ question asks whether every problem whose solutions can be \textit{verified} in polynomial time can also be \textit{solved} in polynomial time. Formally:

\begin{definition}
$\PP$ is the class of languages decidable in polynomial time. $\NP$ is the class of languages for which membership has a polynomial-time verifiable witness~\cite{Arora2009}. $\PP = \NP$ iff every language in $\NP$ is also in $\PP$. By the Cook-Levin theorem~\cite{Cook1971}, this is equivalent to the satisfiability problem (SAT) being solvable in polynomial time.
\end{definition}

Despite decades of effort, the problem remains open. This paper argues that the difficulty is not technical but \textit{conceptual}: previous approaches attempt to prove lower bounds via \textit{constructive} methods (exhibiting specific hard instances, explicit circuit separations, or algebraic obstructions), which are systematically blocked by three barriers. We propose an alternative: reformulate $\PP \neq \NP$ via a \textit{structural entropy argument} that identifies the impossibility of witness compression as the fundamental obstruction, and pinpoint the precise remaining gap that any proof must close.

\subsection{The three barriers}

Three meta-theorems limit the scope of known proof techniques:

\textbf{1. Relativization} \cite{BakerGillSolovay1975}. There exist oracles $A, B$ with $\PP^A = \NP^A$ and $\PP^B \neq \NP^B$. Therefore, any proof of $\PP \neq \NP$ must use techniques that do not relativize -- it must exploit the internal structure of computation, not just input-output behavior.

\textbf{2. Natural proofs} \cite{RazborovRudich1997}. If strong pseudorandom generators exist (a widely believed assumption), then no ``natural'' proof method -- one that is both \textit{constructive} (recognizes hard functions efficiently) and \textit{large} (applies to a dense set of functions) -- can prove superpolynomial circuit lower bounds against $\PP/\poly$.

\textbf{3. Algebrization} \cite{AaronsonWigderson2009}. Even non-relativizing techniques based on arithmetization (as in $\mathrm{IP} = \mathrm{PSPACE}$) cannot separate $\PP$ from $\NP$, because there exist algebraic extensions of oracles in which both $\PP = \NP$ and $\PP \neq \NP$ hold.

\subsection{The positivity alternative}

We observe a \textit{superficial structural parallel} across several major mathematical problems: breakthroughs often succeed not by constructing the desired object, but by showing its non-existence would violate a \textit{positivity condition}. The analogy is at the level of proof strategy (non-constructive arguments via non-negative quantities), not at the level of mathematical content -- the positivity structures involved are fundamentally different in nature:

\begin{center}
\resizebox{\columnwidth}{!}{%
\begin{tabular}{@{}lll@{}}
\hline
Problem & Positive structure & Status \\
\hline
Weil conjectures & Frobenius eigenvalues & Proven \\
BSD (rank $\leq 1$) & N\'eron-Tate height & Proven \\
Hodge conjecture & Hodge-Riemann form & Open \\
RH & Spectral zeta flow & \cite{GeigerRH2026c} \\
\textbf{P vs NP} & \textbf{Kolmogorov complexity} & \textbf{This paper} \\
Broader program & Free energy framework & \cite{FST-Unified2026} \\
\hline
\end{tabular}
}
\end{center}

\noindent\textit{Caveat:} This table illustrates structural parallels at the level of proof strategy; it does not claim that P vs NP is solved or simplified by the positivity perspective. The positivity structures involved are fundamentally different in each case.

The organizing observation: \textit{Kolmogorov complexity serves as a natural non-negative quantity for computational complexity}. It is non-negative ($K(x) \geq 0$), machine-invariant (up to additive constants), and fundamentally non-constructive (not computable) -- properties that allow unbounded $K$ to bypass the three barriers. The resource-bounded variant $K^t$, which is essential for the core results, inherits some but not all of these properties; its barrier status is discussed in Section~\ref{sec:entropy}.

This paper is a \textbf{presentation and framework paper}: it presents the known equivalence between $\PP \neq \NP$ and an entropic condition in a unified framework, but does not prove the separation itself and does not claim the equivalence as a new result (see Related Work, Section~\ref{sec:related}, for the priority of Fortnow~\cite{Fortnow2000} and Orponen et al.~\cite{OrponenKoSchoeningWatanabe1994}). The paper is self-contained; all main results depend only on definitions and proofs within this document. Connections to the broader FST programme are cited for context but are not logically required.


% ===================================================================
% 2. BARRIER ANALYSIS
% ===================================================================
\section{Systematic Barrier Analysis}
\label{sec:barriers}

\subsection{Why constructive methods fail}

All three barriers share a common cause: they block proofs that treat computational problems as \textit{black boxes} or that rely on \textit{statistical} properties of Boolean functions.

\begin{itemize}
\item \textbf{Relativization} blocks proofs that treat the computation as input-output behavior (oracle access), ignoring internal structure.
\item \textbf{Natural proofs} block proofs that identify hard functions via efficiently testable, statistically dense properties.
\item \textbf{Algebrization} blocks proofs that use only the polynomial structure of arithmetized computations.
\end{itemize}

The common thread: all three barriers apply to \textit{model-dependent} arguments -- arguments that work within a specific computational model (oracle machines, circuits, algebraic extensions) rather than targeting the \textit{intrinsic information content} of the computation.

\subsection{The specification for a valid proof}

From the barrier analysis, we derive necessary conditions for any proof of $\PP \neq \NP$:

\begin{enumerate}
\item[\textbf{(S1)}] \textbf{Non-relativizing:} The argument must exploit internal computational structure, not just oracle behavior.
\item[\textbf{(S2)}] \textbf{Non-natural:} The argument must not define a ``largeness'' property that is efficiently computable. It must target individual objects or non-constructive properties.
\item[\textbf{(S3)}] \textbf{Non-algebrizing:} The argument must go beyond polynomial identity testing and arithmetization.
\item[\textbf{(S4)}] \textbf{Model-independent:} The argument must not depend on a specific computational model (circuits, Turing machines, etc.) but on an invariant structural property.
\item[\textbf{(S5)}] \textbf{Robust:} The argument must not break when moving from restricted models (monotone circuits, $\mathrm{AC}^0$) to general computation.
\end{enumerate}

We now show that algorithmic entropy satisfies all five specifications.


% ===================================================================
% 3. ALGORITHMIC ENTROPY
% ===================================================================
\section{Algorithmic Entropy as Positivity Structure}
\label{sec:entropy}

\subsection{Kolmogorov complexity}

Let $U$ be a fixed universal prefix-free Turing machine. The \textit{Kolmogorov complexity} (algorithmic entropy) of a finite binary string $x$ is:
\begin{equation}
K(x) := \min\{|p| : U(p) = x\}\,.
\label{eq:kolmogorov}
\end{equation}

The \textit{conditional} Kolmogorov complexity of $x$ given $y$ is:
\begin{equation}
K(x \mid y) := \min\{|p| : U(p, y) = x\}\,.
\label{eq:conditional_K}
\end{equation}

\begin{theorem}[Invariance {\cite{Kolmogorov1965,Solomonoff1964,Chaitin1966}}]
\label{thm:invariance}
For any two universal prefix-free machines $U, U'$:
\begin{equation}
|K_U(x) - K_{U'}(x)| \leq c_{U,U'}\,,
\end{equation}
where $c_{U,U'}$ is a constant independent of $x$.
\end{theorem}

\begin{theorem}[Incomputability]
\label{thm:incomputable}
$K(x)$ is not computable. More precisely: for every total computable function $f$, there exists $x$ with $K(x) > f(x)$.
\end{theorem}

\begin{theorem}[Incompressibility]
\label{thm:incompressible}
For every $n$, at least $2^n - 2^{n-c+1} + 1$ strings of length $n$ satisfy $K(x) \geq n - c$.
\end{theorem}

\begin{theorem}[Symmetry of Information {\cite{LiVitanyi2008}}]
\label{thm:symmetry}
For all strings $x, y$:
\begin{equation}
K(x, y) = K(x) + K(y \mid x, K(x)) + O(\log(K(x, y)))\,.
\end{equation}
This implies that conditional complexity is well-behaved: the information in the pair $(x, y)$ decomposes additively (up to logarithmic terms) into the information in $x$ plus the additional information in $y$ given $x$.
\end{theorem}

\subsection{Why unbounded Kolmogorov complexity bypasses the barriers}

\textbf{Critical scope limitation.} The barrier immunity analysis below applies to \emph{unbounded} Kolmogorov complexity $K$, which is not the quantity used in the core results of this paper. The core results (Theorem~\ref{thm:high_entropy}, the witness entropy gap) use \emph{resource-bounded} $K^t$, for which barrier immunity is weaker and partially open. This section therefore provides \emph{motivation} for the entropy approach but does \emph{not} establish barrier immunity for the paper's actual technical results. The reader should consult Remark~\ref{rem:barrier-K-vs-Kt} for the precise status of $K^t$.

\begin{claim}[Barrier immunity]
\label{claim:barrier_immune}
Arguments based on \emph{unbounded} Kolmogorov complexity $K$ plausibly satisfy specifications (S1)--(S5). We present heuristic arguments for each; a fully rigorous treatment is an open research question (see Limitations, Section~\ref{sec:discussion}). \textbf{Important caveat:} The core technical results of this paper use \emph{resource-bounded} $K^t$, for which barrier immunity is weaker and partially open -- see Remark~\ref{rem:barrier-K-vs-Kt} below.
\end{claim}

\begin{enumerate}
\item[\textbf{(S1)}] \textbf{Non-relativizing.} $K(x)$ is defined without reference to any oracle -- it is an absolute, fixed quantity for each string $x$. The witness entropy gap (Theorem~\ref{thm:high_entropy}) uses $K^t(w \mid x)$ (resource-bounded conditional complexity). We note an important subtlety: the invariance theorem for resource-bounded $K^t$ is weaker than for unbounded $K$ -- changing the universal machine can introduce a \textit{polynomial} slowdown factor rather than just an additive constant~\cite{LiVitanyi2008}. However, the argument remains non-relativizing because $K^t$ is defined over a fixed computational model without oracle access. A relativizing proof would need to hold for all oracles $A$, but our argument targets oracle-independent quantities: the gap between $O(\log n)$ and $\Omega(|w|)$ is superpolynomial and thus robust under polynomial slowdown factors.

\item[\textbf{(S2)}] \textbf{Non-natural.} $K(x)$ is not computable (Theorem~\ref{thm:incomputable}), so no efficient procedure can test whether a specific function has ``high Kolmogorov complexity.'' This directly violates the ``constructiveness'' requirement of natural proofs.

\item[\textbf{(S3)}] \textbf{Non-algebrizing.} Algebrization applies to proof techniques that use only the polynomial structure of arithmetized computations. Arguments based on $K^t$ operate over the full power of Turing machine computation, not polynomial identity testing. While it remains a research question whether $K^t$-based arguments are formally non-algebrizing in the technical sense of~\cite{AaronsonWigderson2009}, the entropy framework is not naturally expressible as an algebraic identity, suggesting it escapes this barrier.

\item[\textbf{(S4)}] \textbf{Model-independent.} The invariance theorem ensures that $K(x)$ is the same (up to $O(1)$) for all universal machines -- Turing machines, circuits, $\lambda$-calculus, etc.

\item[\textbf{(S5)}] \textbf{Robust.} $K(x)$ does not depend on the computational model used to define it.
\end{enumerate}

\begin{remark}[Barrier immunity: unbounded $K$ vs.\ resource-bounded $K^t$]
\label{rem:barrier-K-vs-Kt}
The barrier immunity claimed above applies to unbounded Kolmogorov complexity $K$, which is an absolute, uncomputable quantity invariant up to $O(1)$ under change of universal machine. However, the core technical results of this paper (Theorem~\ref{thm:high_entropy}, the witness entropy gap) necessarily use \emph{resource-bounded} $K^t$, which is computable and whose invariance is weaker (polynomial slowdown rather than additive constant). The barrier status of $K^t$-based arguments is therefore less clear-cut:
\begin{itemize}
\item \textbf{Relativization:} $K^t(w \mid x)$ is defined without oracles, so our arguments are non-relativizing in the sense that they do not hold ``for all oracles.'' However, $K^{t,A}(w \mid x)$ (with oracle access) can differ substantially from $K^t(w \mid x)$, and a full barrier-immunity proof for the resource-bounded case remains open.
\item \textbf{Natural proofs:} $K^t$ is computable, so unlike $K$, it does not automatically escape the constructiveness requirement. The argument that our approach avoids natural proofs rests on the fact that we do not use $K^t$ as an efficiently testable ``largeness'' property of Boolean functions, but this distinction deserves further formalization.
\item \textbf{Algebrization:} No formal proof exists that $K^t$-based arguments are non-algebrizing in the technical sense of~\cite{AaronsonWigderson2009}. S3 (non-algebrization for $K^t$) remains an open question; we make no claim that our framework provably evades this barrier.
\end{itemize}
In summary: the barrier bypass is rigorous for unbounded $K$ but remains heuristic for resource-bounded $K^t$. Since the witness entropy gap requires $K^t$, establishing formal barrier immunity for the resource-bounded case is an important open problem.
\end{remark}


% ===================================================================
% 4. THE WITNESS ENTROPY GAP
% ===================================================================
\section{The Witness Entropy Gap}
\label{sec:gap}

\subsection{NP witnesses as compressed information}

For an NP language $L$, every yes-instance $x \in L$ has a polynomial-length witness $w$ verifiable in polynomial time:
\begin{equation}
x \in L \iff \exists w,\, |w| \leq \poly(|x|),\, V(x, w) = 1\,.
\end{equation}

The conditional complexity of the witness given the instance measures how much \textit{additional} information the witness contains beyond the instance:
\begin{equation}
\begin{aligned}
K(w \mid x)
&= \text{extra information needed}\\
&\quad \text{to produce } w \text{ from } x\,.
\end{aligned}
\end{equation}

\subsection{What P = NP would imply}

\begin{proposition}[Witness compression under P = NP]
\label{prop:compression}
If $\PP = \NP$, then for every NP language $L$ and every $x \in L$, there exists a polynomial-time canonical witness selector $A$ that produces a valid witness $w$ from $x$:
\begin{equation}
w = A(x)\,.
\end{equation}
Therefore, for some polynomial $q$ depending on the running time of $A$:
\begin{equation}
K^{q(|x|)}(w \mid x) \leq K(A) + O(\log|x|) = O(\log|x|)\,,
\label{eq:compression}
\end{equation}
since the algorithm $A$ is a fixed finite program running in polynomial time, and $O(\log|x|)$ bits encode the input length.
\end{proposition}

\begin{proof}
If $\PP = \NP$, the search version of SAT is solvable in polynomial time (by self-reducibility). Given a satisfiable formula $\varphi$, the algorithm $A$ finds a satisfying assignment bit-by-bit in polynomial time. The program ``run $A$ on input $x$'' has length $|A| + O(\log|x|)$, where $|A|$ is a constant.
\end{proof}

\subsection{The entropy obstruction}

\begin{remark}[Unbounded vs.\ resource-bounded complexity -- why $K^t$ is essential]
\label{rem:bounded-vs-unbounded}
The distinction between $K(w \mid x)$ and $K^t(w \mid x)$ is crucial and must be understood before the main theorem. For any decidable language $L$ and any $(x, w)$ with $V(x,w) = 1$, the brute-force verifier-search has $K(w \mid x) \leq O(1)$ (the search program has constant length -- it simply enumerates all candidates and checks each). The information-theoretic content of NP-hardness lies entirely in the \textit{time} required to extract the witness, not in the \textit{description length}. This is why $K^t$ with polynomial $t$ is the correct measure: it captures the computational cost of witness extraction, which is the essence of the $\PP$ vs $\NP$ question.
\end{remark}

\begin{theorem}[Superlogarithmic witness obstruction -- resource-bounded, conditional {\normalfont (cf.\ Fortnow~\cite{Fortnow2000}, Allender~\cite{Allender2001})}]
\label{thm:high_entropy}
The following is a reformulation of known connections between resource-bounded Kolmogorov complexity and P vs NP (see Related Work, Section~\ref{sec:related}) in the specific language of conditional witness entropy.

If $\PP \neq \NP$, then for SAT, for every polynomial $q$ and every constant $C>0$, there exist infinitely many satisfiable formulas $x$ such that \textbf{every} satisfying assignment $w$ satisfies:
\begin{equation}
K^{q(|x|)}(w \mid x) > C \log |x|\,.
\label{eq:high_entropy}
\end{equation}
That is, for every fixed polynomial time bound and every fixed logarithmic advice constant, infinitely many instances have no valid witness with that small a time-bounded conditional description. The same statement holds for any NP-complete search problem to which SAT reduces by a polynomial-time witness-preserving search reduction.

Equivalently (by contraposition): if there exist a polynomial $q$ and a constant $C$ such that every satisfiable formula $x$ has some satisfying assignment $w$ with $K^{q(|x|)}(w \mid x) \leq C\log |x|$, then $\PP = \NP$.

Note on quantifier order: this theorem is deliberately stated for each fixed pair $(q,C)$. It does \emph{not} assert a single infinite instance family that resists all polynomial time bounds simultaneously, and it does \emph{not} assert near-maximal entropy of witnesses.
\end{theorem}

\begin{proof}
We prove the contrapositive for SAT. Suppose some polynomial $q$ and constant $C$ have the property that every satisfiable formula $x$ has a satisfying assignment $w$ with $K^{q(|x|)}(w \mid x) \leq C\log |x|$. Then for every satisfiable $x$, there exists a program $p_x$ of length at most $C\log |x|$ that produces a satisfying assignment within $q(|x|)$ steps. There are at most $|x|^C$ such programs. A polynomial-time SAT search algorithm can enumerate all programs of length at most $C\log |x|$, run each for $q(|x|)$ steps on input $x$, and verify each output. Since SAT search is polynomial-time equivalent to SAT decision by self-reducibility, this would imply $\PP=\NP$.

The ``infinitely many'' clause follows by the same contrapositive plus finite hardcoding: if only finitely many formulas failed the bound, their witnesses could be stored in a finite lookup table and the above enumeration would solve all remaining instances in polynomial time. Constructing such hard instances explicitly is not provided by the argument; the theorem is a non-constructive reformulation.

\textit{Why unbounded $K$ fails:} For any SAT instance $x$ with a satisfying assignment $w$, the brute-force program ``enumerate all assignments, test each, output the first satisfying one'' has length $O(1)$ and finds $w$ from $x$ in time $O(2^{|x|})$. Thus $K(w \mid x) \leq O(1)$ -- the unbounded conditional complexity is trivially small. This is why the resource bound is essential: the information-theoretic content of NP-hardness lies in the \textit{time} required to extract the witness, not in the description length.
\end{proof}

\begin{conjecture}[Near-maximal witness entropy]
\label{conj:near-maximal-witness-entropy}
For some NP-complete search formulation, and for every polynomial $q$, there exist infinitely many yes-instances $x$ such that every valid witness $w$ satisfies
\begin{equation}
K^{q(|x|)}(w \mid x) \geq |w| - O(\log |x|)\,.
\end{equation}
This near-maximal strengthening would yield the sharp $\Omega(|w|)$ vs.\ $O(\log |x|)$ entropy gap used heuristically in several bridge targets below. It is not implied by Theorem~\ref{thm:high_entropy}; it is an additional open hypothesis.
\end{conjecture}


\subsection{The gap}

Combining Proposition~\ref{prop:compression} and Theorem~\ref{thm:high_entropy}:

\begin{corollary}[Witness Entropy Gap -- the core contradiction structure]
\label{cor:gap}
For any fixed polynomial time bound $q$ and logarithmic constant $C$, Theorem~\ref{thm:high_entropy} gives infinitely many SAT instances whose valid witnesses are not $C\log |x|$-compressible within time $q(|x|)$. For unique-witness instances, the contradiction structure is sharpest:
\begin{itemize}
\item $\PP = \NP$ $\Rightarrow$ a polynomial-time canonical witness selector $A$ exists with $K^{q(|x|)}(A(x) \mid x) \leq O(\log |x|)$ (Proposition~\ref{prop:compression}).
\item $\PP \neq \NP$ $\Rightarrow$ for every fixed $(q,C)$, infinitely many instances have no valid witness $w$ with $K^{q(|x|)}(w \mid x) \leq C\log |x|$ (Theorem~\ref{thm:high_entropy}).
\end{itemize}
This is not a proof of $\PP \neq \NP$, but a characterization of the logarithmic-compression threshold: $\PP = \NP$ if and only if a polynomial-time canonical witness selector with logarithmic conditional description exists for SAT. The stronger $\Omega(|w|)$ vs.\ $O(\log n)$ gap follows only under the near-maximal witness entropy hypothesis (Conjecture~\ref{conj:near-maximal-witness-entropy}).

The gap is sharpest for instances with a unique satisfying assignment; for instances with multiple witnesses, the theorem rules out \emph{any} logarithmically short polynomial-time witness description on the hard instances.
\end{corollary}

\begin{remark}[Unique vs.\ multiple witnesses]
Theorem~\ref{thm:high_entropy} guarantees instances where \textit{every} valid witness has superlogarithmic conditional complexity relative to the fixed time bound and constant under consideration. For these instances, the contradiction structure holds regardless of how many witnesses exist: any algorithm $A$ must output \textit{some} valid witness, and none of the valid witnesses has the prohibited logarithmically short description.

For unique-witness instances, the structure is especially clean: $A$ has no choice over which witness to produce. In general, however, the high-entropy instances guaranteed by Theorem~\ref{thm:high_entropy} exist non-constructively (via contraposition), and it remains open whether they can be constructed explicitly -- see Section~\ref{sec:obstruction}.
\end{remark}


% ===================================================================
% 5. ALGORITHMIC POSITIVITY
% ===================================================================
\section{Algorithmic Positivity}
\label{sec:positivity}

We now formalize the No-Go framework.

\subsection{Definition}

\begin{definition}[Algorithmic Positivity]
\label{def:alg_pos}
An NP language $L$ (with verifier $V$) is \textit{algorithmically positive} if there exists a \textit{witness-producing} function $W: \{0,1\}^* \to \{0,1\}^*$ satisfying (AP1) and (AP2):

\begin{enumerate}
\item[\textbf{(AP1)}] \textbf{Polynomial scaling.} $W$ is computable in time $\poly(|x|)$, and for every $x \in L$: $V(x, W(x)) = 1$.

\item[\textbf{(AP2)}] \textbf{Entropic convexity.} There exists a polynomial $q$ such that for every $x \in L$:
\begin{equation}
K^{q(|x|)}(W(x) \mid x) \leq O(\log|x|)\,.
\end{equation}
The witness production is ``entropically cheap'': the witness $W(x)$ adds at most logarithmic resource-bounded information beyond the input. (Note: since $W$ is a fixed polynomial-time program, the polynomial $q$ is determined by $W$ and does not depend on the instance $x$.)
\end{enumerate}
\end{definition}

\begin{remark}
Note that (AP2) targets the \textit{witness} $W(x)$, not just the decision bit. A single decision bit $D(x) \in \{0,1\}$ trivially satisfies $K(D(x) \mid x) \leq 1$, so entropic convexity would be vacuous for decision functions. The meaningful measure is the resource-bounded complexity of producing a \textit{witness} -- this is where the information-theoretic content of NP-hardness resides.
\end{remark}

\begin{remark}
Every $\PP$-decidable NP language is algorithmically positive: if $L \in \PP \cap \NP$, then the polynomial-time search algorithm (obtained via self-reducibility for NP-complete problems, or directly for natural NP problems) produces a witness $W(x)$ with $K^{q(|x|)}(W(x) \mid x) \leq |W_{\text{prog}}| + O(\log|x|) = O(\log|x|)$.
\end{remark}

\begin{remark}[Heuristic properties of P-decidable NP languages]
\label{rem:heuristic-ap}
For NP languages in $\PP$, the witness-producing function $W$ additionally satisfies two heuristic properties that capture the ``structural simplicity'' of polynomial-time witness extraction:

\textbf{(H1) Compositional stability.} For structured instances under a well-defined composition operator $\circ$, the witness for the composed input is typically determined by the component witnesses with at most logarithmic overhead: $K^{\poly}(W(x_1 \circ x_2) \mid W(x_1), W(x_2)) \leq O(\log(|x_1|+|x_2|))$.

\textbf{(H2) Monotone stability.} Small perturbations of the input produce witnesses of bounded additional complexity: $d_H(x,x') \leq 1 \implies K(W(x) \mid W(x'), x) \leq O(\mathrm{poly}(\log|x|))$.

These properties are motivated by the structure of $\NP$-complete problems but are not used in the formal No-Go argument (Theorem~\ref{thm:nogo}). They serve as additional heuristic evidence for the Entropic Separation Conjecture.
\end{remark}

\subsection{The Entropic Separation Conjecture}

\begin{conjecture}[Entropic Separation -- ESC]
\label{conj:entropic}
No NP-complete language is algorithmically positive. That is: for every NP-complete language $L$ and every polynomial-time witness-producing function $W$, there exist infinitely many yes-instances $x \in L$ such that $W$ fails to produce a valid witness for $x$ (i.e., $V(x, W(x)) = 0$), or $K^{q(|x|)}(W(x) \mid x) \geq \omega(\log|x|)$ for all polynomials $q$.

An equivalent formulation in terms of truth-table complexity (see Theorem~\ref{thm:slice_nogo}): for some NP-complete language $L$, for every polynomial $q$, infinitely many $n$ satisfy $K^{q(2^n)}(L_n) \geq n$. (Both formulations are equivalent to $\PP \neq \NP$; their mutual equivalence follows transitively.)
\end{conjecture}

\begin{proposition}
\label{prop:equiv}
The Entropic Separation Conjecture~\ref{conj:entropic} implies $\PP \neq \NP$.
\end{proposition}

\begin{proof}
If $\PP = \NP$, then for every NP-complete language $L$, the search version is solvable in polynomial time (for SAT, by self-reducibility; for general NP-complete $L$, by reducing $L$ to SAT, applying the SAT search algorithm, and reconstructing the witness). This gives a polynomial-time witness-producing function $W$ satisfying (AP1). Since $W$ is a fixed program running in polynomial time, $K^{q(|x|)}(W(x) \mid x) \leq |W_{\text{prog}}| + O(\log|x|) = O(\log|x|)$, satisfying (AP2). Therefore, NP-complete languages would be algorithmically positive, contradicting the conjecture.
\end{proof}

\begin{remark}[ESC is equivalent to $\PP \neq \NP$]
\label{rem:esc-equivalence}
The converse of Proposition~\ref{prop:equiv} also holds: $\PP \neq \NP$ implies ESC. Indeed, if $\PP = \NP$, then every NP-complete language $L$ has a polynomial-time algorithm, giving $K^{\mathrm{poly}(2^n)}(L_n) = O(\log n)$ by Lemma~\ref{lem:p_low_entropy}, which violates the ESC condition $K^{q(2^n)}(L_n) \geq n$. Thus ESC $\Leftrightarrow$ $\PP \neq \NP$: the Entropic Separation Conjecture is a \textit{reformulation} of the P vs NP problem in information-theoretic language, not an independent strengthening.

\textit{Priority note:} This equivalence is a reformulation of the known equivalence between P $\neq$ NP and unbounded time-bounded instance complexity of SAT, established by Fortnow~\cite{Fortnow2000} building on the instance complexity framework of Orponen et al.~\cite{OrponenKoSchoeningWatanabe1994}. Our contribution is the specific formulation in terms of conditional witness entropy $K^t(w \mid x)$ rather than instance complexity $\mathrm{ic}_t(x : \mathrm{SAT})$.
\end{remark}

\subsection{Why NP-complete problems resist algorithmic positivity}

The failure of each condition for NP-complete problems:

\textbf{(AP2) fails:} For NP-complete problems, producing a witness $w$ for $x \in L$ requires extracting information that is not uniformly logarithmically compressible relative to the instance. Theorem~\ref{thm:high_entropy} shows, conditional on $\PP \neq \NP$, that for every fixed polynomial time bound and logarithmic constant, infinitely many SAT instances have no valid witness within that description budget. The stronger near-maximal bound $K^t(w \mid x)=\Omega(|w|)$ is an additional bridge hypothesis, not a theorem. No polynomial-time witness-producing function $W$ can achieve $K^{\poly}(W(x) \mid x) \leq O(\log|x|)$ for all SAT instances.

As additional heuristic evidence (not part of the formal argument), the properties (H1)--(H2) from Remark~\ref{rem:heuristic-ap} also fail for NP-complete problems:

\textbf{(H1) fails:} SAT is not compositionally decomposable in the required sense. The satisfiability of a conjunction $\varphi_1 \wedge \varphi_2$ is not determined by the satisfiability of $\varphi_1$ and $\varphi_2$ individually (they share variables). This non-decomposability is precisely the source of NP-hardness.

\textbf{(H2) fails for random $k$-SAT near the threshold:} For random $k$-SAT instances near the satisfiability threshold, SAT exhibits extreme sensitivity in its search version: flipping a single clause can restructure the entire witness space. Specifically, near the phase transition (clause-to-variable ratio $\alpha \approx \alpha_c(k)$), a single clause addition can fragment the solution cluster into exponentially many components, producing witnesses of unbounded additional complexity~\cite{AchlioptasRicci2006,MezardZecchina2002}. \textbf{Scope limitation:} This instability is rigorously established only for random $k$-SAT in the near-threshold regime $\alpha \approx \alpha_c(k)$. For structured, planted, or under-constrained instances ($\alpha \ll \alpha_c$), the witness space may be well-connected, and monotone stability can hold. The heuristic argument for (H2) failure as a general NP-hardness phenomenon thus rests on the assumption that the near-threshold regime is representative of the combinatorial difficulty of NP-complete problems.


% ===================================================================
% 6. THE NO-GO THEOREM
% ===================================================================
\section{The No-Go Theorem}
\label{sec:nogo}

\begin{theorem}[Entropic No-Go -- Reformulation]
\label{thm:nogo}
The following are equivalent:
\begin{enumerate}
\item[(i)] $\PP \neq \NP$.
\item[(ii)] The Entropic Separation Conjecture~\ref{conj:entropic} holds.
\item[(iii)] No NP-complete language is decidable in polynomial time. That is: for every NP-complete language $L$ and every polynomial-time algorithm $A$, there exist instances $x$ where $A(x) \neq L(x)$.
\end{enumerate}
The equivalence $(i) \Leftrightarrow (ii)$ shows that $\PP \neq \NP$ is \textit{equivalent} to an entropy obstruction: polynomial-time algorithms cannot compress NP-witnesses below their intrinsic Kolmogorov complexity.
\end{theorem}

\begin{proof}
$(ii) \Rightarrow (i)$: Shown in Proposition~\ref{prop:equiv}. Contrapositive: if $\PP = \NP$, then every NP-complete language has a polynomial-time witness-producing function (via the SAT search algorithm and polynomial-time reductions), satisfying both (AP1) and (AP2). This makes NP-complete languages algorithmically positive, contradicting ESC.

$(i) \Rightarrow (ii)$: Contrapositive: suppose ESC fails, i.e., some NP-complete language $L$ is algorithmically positive. Then $L$ has a polynomial-time witness-producing function $W$ with $K^{\poly}(W(x) \mid x) \leq O(\log|x|)$. In particular, $W$ is a polynomial-time search algorithm for $L$. Since $L$ is NP-complete, this implies $\PP = \NP$.

$(i) \Leftrightarrow (iii)$: (iii) states that for every NP-complete language $L$, $L \notin \PP$. This is equivalent to $\PP \neq \NP$ by the existence of NP-complete languages (if any NP-complete $L$ were in $\PP$, then $\PP = \NP$; conversely, $\PP = \NP$ would place every NP-complete language in $\PP$).

The content of the theorem is not a new separation result, but a \textit{unified presentation} of the known equivalence between $\PP \neq \NP$ and entropic incompressibility of NP-witnesses (cf.\ Fortnow~\cite{Fortnow2000}, Orponen et al.~\cite{OrponenKoSchoeningWatanabe1994}). The specific framing in terms of algorithmic positivity and the No-Go structure is the contribution of this paper.
\end{proof}

\subsection{Exclusion of alternative scenarios}

\begin{proposition}[Excluded scenarios]
\label{prop:excluded}
The following scenarios are incompatible with the conjunction of entropic positivity and Kolmogorov obstruction:

\begin{enumerate}
\item[\textbf{(X1)}] \textbf{Polynomial algorithm with exponential entropy production.} An algorithm that ``invents'' new information during computation. Excluded because for any deterministic algorithm $A$ running in time $t_A(|x|)$, the resource-bounded conditional complexity satisfies $K^{t_A(|x|)}(A(x) \mid x) \leq |A| + O(\log|x|)$: the program ``run $A$ on input $x$'' has constant description length. In particular, the \textit{conditional} algorithmic information of $A(x)$ given $x$ is at most $O(1)$ bits at the time scale of $A$'s own computation -- meaning $A$ cannot ``create'' information that was not already implicit in the input. (For randomized algorithms, the random bits $r$ provide additional input, so $K(A(x,r) \mid x)$ can be high. However, the P vs NP question concerns deterministic computation; moreover, under widely believed derandomization assumptions, $\BPP = \PP$~\cite{ImpagliazzoWigderson1997}, so randomness does not provide additional power.)

\item[\textbf{(X2)}] \textbf{Compositional collapse of NP.} All SAT instances decompose into independently solvable parts. Excluded by the NP-completeness of SAT: if SAT decomposed compositionally, then all NP problems would decompose (via reduction). However, general SAT instances involve shared variables across clauses, so the satisfiability of a conjunction $\varphi_1 \wedge \varphi_2$ is not determined by the satisfiability of the parts individually. The Valiant-Vazirani theorem~\cite{ValiantVazirani1986} further strengthens this: SAT reduces (under randomized reductions) to instances with a unique satisfying assignment, where compositional decomposition is impossible since the unique witness encodes global correlations across all variables.

\item[\textbf{(X3)}] \textbf{Smooth decision landscape.} All SAT instances have monotonically stable decision boundaries. Excluded by phase-transition phenomena in random SAT: near the satisfiability threshold, the decision changes discontinuously under small clause additions.
\end{enumerate}
\end{proposition}


% ===================================================================
% 7. RESOURCE-BOUNDED SLICE ENTROPY
% ===================================================================
\section{Resource-Bounded Slice Entropy}
\label{sec:slices}

The witness entropy gap of Section~\ref{sec:gap} targets individual instances. A cleaner formulation targets the \textit{truth table} of the language on all inputs of a given length.

\subsection{Slices of a language}

\begin{definition}[Slice]
For a language $L \subseteq \{0,1\}^*$, the $n$-th \textit{slice} $L_n \in \{0,1\}^{2^n}$ is the truth table of $L$ on inputs of length $n$: $(L_n)_i = 1 \iff x_i \in L$, where $x_1, \ldots, x_{2^n}$ are the $n$-bit strings in lexicographic order.
\end{definition}

\begin{definition}[Resource-bounded algorithmic entropy]
For a time bound $t(\cdot)$:
\begin{equation}
K^t(x) := \min\{|p| : U(p) = x \text{ in at most } t(|x|) \text{ steps}\}\,.
\label{eq:Kt}
\end{equation}
\end{definition}

\subsection{The slice entropy lemma}

\begin{lemma}[P implies low slice entropy]
\label{lem:p_low_entropy}
If $L \in \PP$, then there exists a polynomial $q$ and a constant $c$ such that for all $n$:
\begin{equation}
K^{q(2^n)}(L_n) \leq c + \lceil \log n \rceil\,.
\label{eq:p_entropy}
\end{equation}
\end{lemma}

\begin{proof}
Let $A$ be a polynomial-time algorithm deciding $L$ in time $n^d$. The program $p$: ``read $n$ from input; for each $x \in \{0,1\}^n$, run $A(x)$ and output the result bit'' has length $|A| + O(\log n)$ (to encode $n$). Its runtime is $2^n \cdot n^d = \poly(2^n)$. Thus $K^{q(2^n)}(L_n) \leq |A| + O(\log n) = c + \lceil \log n \rceil$.
\end{proof}

\subsection{The No-Go theorem via slice entropy}

\begin{theorem}[Slice Entropy No-Go]
\label{thm:slice_nogo}
If there exists an $\NP$-complete language $L$ satisfying the \textit{Entropy Hardness Hypothesis}:

\textbf{(EH)} For every polynomial $q$, there exist infinitely many $n$ with
\begin{equation}
K^{q(2^n)}(L_n) \geq n\,,
\label{eq:eh}
\end{equation}

then $\PP \neq \NP$.
\end{theorem}

\begin{proof}
Immediate from Lemma~\ref{lem:p_low_entropy}: $L \in \PP$ would give $K^{\poly(2^n)}(L_n) = O(\log n)$, contradicting~\eqref{eq:eh} for large $n$.
\end{proof}

\begin{remark}
The Entropy Hardness Hypothesis (EH) is a concrete, well-defined conjecture about a specific NP-complete language (e.g., SAT). It asserts that the truth table of SAT on $n$-bit formulas cannot be compressed to fewer than $n$ bits, even when the compression algorithm is allowed $\poly(2^n)$ time.

This is the \textit{cleanest} No-Go formulation: the ``easy direction'' ($L \in \PP \Rightarrow$ low entropy) is a theorem; the ``hard direction'' (EH for SAT) is the entire content of $\PP \neq \NP$.

Note that the threshold $n$ in EH is much smaller than the truth table length $2^n$: a random string of length $2^n$ has $K \approx 2^n$. The threshold $n$ is chosen because Lemma~\ref{lem:p_low_entropy} gives $K^{\poly(2^n)}(L_n) = O(\log n)$ for $L \in \PP$; any threshold $\omega(\log n)$ would suffice to separate from P, and $n$ provides a clean, moderate lower bound that is already superpolynomial in the input length.
\end{remark}

\begin{remark}[Encoding sensitivity of EH]
\label{rem:encoding}
The Entropy Hardness Hypothesis depends on the encoding of SAT instances as binary strings. Different encodings yield truth tables $\SAT_n$ of different lengths and potentially different Kolmogorov complexities. The hypothesis should be understood as: there exists a \emph{standard} polynomial-time encoding $\mathrm{enc}: \mathrm{Formulas} \to \{0,1\}^*$ such that $K^{q(2^n)}(\SAT_n^{\mathrm{enc}}) \geq n$ for infinitely many $n$. For resource-bounded $K^t$, changing the universal machine introduces a polynomial slowdown rather than just an additive constant. However, since EH quantifies over \emph{all} polynomials $q$, a polynomial slowdown in the time bound does not affect the hypothesis: if $K^{q(2^n)}(\SAT_n) \geq n$ for all $q$, the same holds after re-parameterization. The \emph{length} of the truth table (and hence the threshold $n$) is encoding-dependent.
\end{remark}

\begin{remark}[Connection to circuit lower bounds]
\label{rem:circuit-connection}
The Entropy Hardness Hypothesis is closely related to circuit lower bounds. A Boolean function $f_n$ computed by a circuit of size $s$ has truth table describable in $O(s \log s)$ bits, so $K^{O(s \cdot 2^n)}(f_n) \leq O(s \log s) + O(\log n)$. Thus:
\begin{itemize}
\item EH with threshold $n$ implies $\SAT_n$ requires circuits of size $s \geq 2^{\Omega(n / \log n)}$, i.e., $\NP \not\subseteq \mathrm{SIZE}(2^{o(n)})$.
\item In particular, EH implies $\NP \not\subseteq \PP/\poly$ (since polynomial-size circuits give $K^{\poly(2^n)}(f_n) \leq O(n^c \log n)$, which is $o(n)$ for large $n$).
\end{itemize}
This shows that the entropy framework and the circuit complexity approach are two perspectives on the same underlying separation.
\end{remark}

\begin{remark}[SAT slice entropy at the phase transition -- exploratory numerics]
\label{rem:sat-entropy}
As a purely illustrative exercise (not constituting evidence for or against the ESC), we computed the Shannon entropy $H_{\mathrm{slice}}$ of the truth table of random 3-SAT at the phase transition $\alpha \approx 4.267$ for $n = 8, 10, 12, 14, 16$. The entropy grows from $H \approx 1.5$ ($n=8$) to $H \approx 2.3$ ($n=14$). \textbf{Limitations:} (i) Only 5 data points are available, which is far too few to determine asymptotic growth behaviour -- any monotone function (including logarithmic growth, which would \emph{refute} the ESC) is consistent with these data. (ii) The Shannon entropy $H_{\mathrm{slice}}$ of the truth table is \emph{not} the same as the Kolmogorov complexity $K^t(\mathrm{SAT}_n)$; the relationship is not formally established here. (iii) $n \leq 16$ is far below the asymptotic regime relevant to the ESC. Full validation would require $n \geq 30$, which is computationally infeasible. Script: \texttt{compute\_sat\_entropy.py}.
\end{remark}

\subsection{Why (EH) should be true}

The truth table $\SAT_n$ has length $2^n$ and describes the satisfiability of all $n$-bit encodings of Boolean formulas. There are $2^{2^n}$ possible truth tables of this length, and $\SAT_n$ is a \textit{specific} one determined by a complete mathematical definition ($\SAT$). Nevertheless:

\begin{enumerate}
\item \textbf{Pseudorandom structure.} The satisfiability boundary in the space of formulas exhibits phase-transition behavior \cite{Monasson1999}: near the threshold, the pattern of satisfiable vs.\ unsatisfiable instances has high apparent randomness.

\item \textbf{No known compression.} Despite decades of SAT-solver research, no method compresses $\SAT_n$ to $o(2^n)$ bits in $\poly(2^n)$ time. The best algorithms (CDCL solvers) run in time $2^{\Theta(n)}$ in the worst case.

\item \textbf{Cryptographic evidence.} If (EH) fails, then efficient PRGs from one-way functions would not exist \cite{ImpagliazzoWigderson1997}, collapsing much of modern cryptography.
\end{enumerate}


% ===================================================================
% 8. THE PRECISE OBSTRUCTION
% ===================================================================
\section{The Precise Obstruction}
\label{sec:obstruction}

\subsection{The uniformity gap}

The main argument (Section~\ref{sec:gap}) has a precise gap that we now identify.

\textbf{What is proven:} There \textit{exist} NP instances whose witnesses have high conditional Kolmogorov complexity (Theorem~\ref{thm:high_entropy}). This is a \textit{non-uniform} statement about individual strings.

\textbf{What is needed:} A \textit{uniform} statement: for every polynomial-time algorithm $A$, there exist infinitely many instances $x$ where $A(x)$ fails. This requires converting the existential entropy bound into an \textit{adversarial} bound against all efficient algorithms simultaneously.

\begin{definition}[The uniformity bridge]
\label{def:uniformity}
Let $L$ be an NP-complete language with verifier $V$. The \textit{uniformity bridge} is the statement: for every polynomial-time computable function $A: \{0,1\}^* \to \{0,1\}^*$, there exist infinitely many yes-instances $x \in L$ such that $A$ fails to produce a valid witness:
\begin{equation}
\forall A \in \mathrm{FP},\, \exists^\infty x \in L:\, V(x, A(x)) = 0\,.
\end{equation}
Equivalently, in entropic terms: for every polynomial-time $A$ and every polynomial $q$, there exist infinitely many $x \in L$ such that every valid witness $w$ for $x$ satisfies $K^{q(|x|)}(w \mid x) > |A| + O(\log|x|)$.

Note: the uniformity bridge is formulated as a search statement ($A$ must produce a \textit{witness}, i.e., $A \in \mathrm{FP}$ and $V(x, A(x)) = 1$). The connection to the decision problem $\PP \neq \NP$ follows from the NP-completeness of $L$: for NP-complete languages, the search and decision versions are polynomial-time equivalent (via self-reducibility for SAT, or more generally via Cook-Levin plus polynomial-time reductions).
\end{definition}

The uniformity bridge is the P-vs-NP analog of:
\begin{itemize}
\item The ``higher Gross-Zagier formula'' in BSD (coupling analytic rank to algebraic rank).
\item The ``hard direction'' in the Hodge conjecture (showing positivity of the Hodge-Riemann form implies algebraicity).
\end{itemize}

In each case, the positivity structure is established, but a \textit{coupling theorem} is needed to convert positivity into a definitive result.

\subsection{Known partial results}

The uniformity bridge is partially established in restricted settings:

\begin{enumerate}
\item \textbf{Resolution proof systems:} Exponential lower bounds on resolution refutations of random $k$-SAT instances near the threshold \cite{BenSasson2001}. This proves that resolution-based SAT solvers cannot solve these instances efficiently.

\item \textbf{Bounded-depth circuits:} $\SAT \notin \mathrm{AC}^0$ (Furst-Saxe-Sipser 1984, Ajtai 1983; see~\cite{Hastad1987} for optimal bounds). This proves the entropy obstruction for constant-depth circuits.

\item \textbf{Monotone circuits:} Exponential lower bounds for monotone circuits computing CLIQUE \cite{Razborov1985}; related bounded-depth lower bounds \cite{Razborov1987}.

\item \textbf{Algebraic proof systems:} Exponential lower bounds for Nullstellensatz and Polynomial Calculus refutations \cite{Grigoriev2001}.
\end{enumerate}

Each of these is a \textit{model-specific} instantiation of the entropy obstruction. The challenge is to prove it for \textit{general} polynomial-time computation -- which is precisely the $\PP \neq \NP$ problem.

\subsection{Why the gap is hard}

The uniformity gap is hard for the same reason that $\PP$ vs $\NP$ is hard: it requires a statement about \textit{all} polynomial-time algorithms. In the uniform model, this is a universal quantifier over a countably infinite, recursively enumerable set (all Turing machines running in polynomial time). In the non-uniform model ($\PP/\poly$), the quantifier ranges over all polynomial-size circuit families, an uncountable set.

The entropy approach \textit{reformulates} the problem but does not \textit{simplify} it:
\begin{quote}
\textit{P vs NP is equivalent to asking whether Kolmogorov complexity can be polynomially compressed for NP-witness structures~\cite{Fortnow2000}. The entropy framework makes this equivalence explicit in the language of witness entropy and identifies concrete bridge strategies, but converting this to a proof against all algorithms remains the central challenge -- which is P vs NP itself in different language.}
\end{quote}

\subsection{Four candidate bridge strategies}

We identify four concrete research directions for closing the uniformity gap:

\textbf{(B1) Dense families with unique random witnesses.} Construct \textit{dense} families $S_n \subseteq \{0,1\}^{\poly(n)}$ with $|S_n|/2^{\poly(n)}$ non-vanishing, such that every $x \in S_n$ has a unique valid witness $w_x$ with $K^{\poly(n)}(w_x \mid x) \geq |w_x| - O(\log n)$. Density ensures that no polynomial-time algorithm can ``avoid'' all hard instances. The construction requires efficient reductions from Kolmogorov-random strings to SAT instances with unique solutions -- a non-trivial but plausible coding-theoretic task.

\textbf{(B2) Self-referential diagonalization.} For each polynomial-time machine $A$, construct an instance $x_A$ whose unique valid witness $w$ contains (or encodes) a description of $A$ itself, so that $K^{\poly(|x_A|)}(w \mid x_A) \geq |A|$. This creates a self-referential obstruction: $A$ cannot produce a witness that is ``about itself'' without exceeding its own description length. The approach combines G\"odel-style diagonalization with Kolmogorov complexity and avoids relativization because the self-reference targets the machine's \textit{description}, not its oracle behavior.

\textbf{(B3) Resource-bounded Kolmogorov lower bounds on slices.} Prove that for SAT (or another NP-complete language), the resource-bounded complexity satisfies $K^{\poly(2^n)}(\SAT_n) \geq 2^{n-o(n)}$ for infinitely many $n$. This is a \textit{uniform} statement about the truth table and directly implies $\PP \neq \NP$ via the Slice Entropy No-Go (Theorem~\ref{thm:slice_nogo}). The strategy connects to proof complexity: if every short proof of an instance's satisfiability has high Kolmogorov complexity, then no polynomial-time algorithm can systematically produce such proofs.

\textbf{(B4) The Minimum Circuit Size Problem (MCSP).} Recent work connects resource-bounded Kolmogorov complexity to circuit complexity via the \textit{Minimum Circuit Size Problem} ($\mathrm{MCSP}$): given a truth table $t \in \{0,1\}^{2^n}$ and a parameter $s$, decide whether there exists a circuit of size $\leq s$ computing $t$. Allender et al.~\cite{Allender2006} established deep connections between $\mathrm{MCSP}$ and $K^t$. While $\mathrm{MCSP} \in \NP$, its complexity status is open: it is not known to be $\NP$-complete, yet proving $\mathrm{MCSP} \in \PP$ would have strong derandomisation consequences. If $\mathrm{MCSP}$ is hard for some complexity class (e.g., $\mathrm{MCSP} \notin \PP/\mathrm{poly}$), this would yield circuit lower bounds that imply the Entropy Hardness Hypothesis. The chain is indirect but concrete: \textit{hardness of MCSP $\Rightarrow$ circuit lower bounds $\Rightarrow$ truth-table incompressibility $\Rightarrow$ EH}.


% ===================================================================
% 9. CONNECTION TO PROOF COMPLEXITY
% ===================================================================
\section{Proof Complexity as a Testing Ground}
\label{sec:proof_complexity}

\subsection{SAT algorithms and short proofs}

A deep connection exists between algorithms and proofs:
\begin{equation}
\begin{aligned}
\SAT \in \PP
\implies{}& \text{every unsatisfiable formula}\\
& \text{has a short refutation,}
\end{aligned}
\end{equation}
in a proof system at least as strong as Extended Frege.

\begin{proposition}[Proof complexity formulation -- Cook {\cite{Cook1975}}]
\label{prop:proof_complexity}
If $\PP \neq \NP$, then there is no polynomially bounded proof system for tautologies. Equivalently: $\PP = \NP$ if and only if there exists a proof system in which every tautology of length $n$ has a proof of length $\poly(n)$.
\end{proposition}

This is Cook's reformulation~\cite{Cook1975}: proving superpolynomial lower bounds on proof length in every propositional proof system is equivalent to showing $\PP \neq \coNP$ (which follows from $\PP \neq \NP$).

\subsection{Entropy of proofs}

The entropy framework connects to proof complexity via the following question: for a random unsatisfiable formula $\varphi$ and its shortest refutation $\pi$ in a given proof system, what is $K(\pi \mid \varphi)$? If $\pi$ is ``algorithmically random'' relative to $\varphi$ -- meaning $K(\pi \mid \varphi)$ is close to $|\pi|$ -- then no efficient algorithm can systematically produce such refutations.

This provides a concrete testing ground: if one can show that for random unsatisfiable $3$-SAT near the threshold, the shortest refutation in Extended Frege has both superpolynomial length and high conditional Kolmogorov complexity, this would be strong evidence for the entropy obstruction.

\begin{conjecture}[Resolution width and Kolmogorov complexity]
\label{conj:resolution-width}
For unsatisfiable CNF formulas $\varphi_n$ on $n$ variables, the minimum resolution width $w(\varphi_n \vdash \bot)$ satisfies
\[
w(\varphi_n \vdash \bot) \geq \Omega\left( K(\pi^* | \varphi_n) / \log n \right),
\]
where $\pi^*$ is the shortest resolution refutation of $\varphi_n$. In particular, formulas whose refutations have high Kolmogorov complexity require wide clauses.
\end{conjecture}

\textit{Motivation.} Ben-Sasson and Wigderson~\cite{BenSasson2001} showed that resolution width lower bounds imply resolution length lower bounds via $|\pi| \geq 2^{w - O(n)}$. Our conjecture reverses the direction: high complexity of the proof object (measured by $K$) forces width. If true, this would provide an information-theoretic characterisation of resolution width, connecting proof complexity to the ESC framework.

\begin{remark}[Resolution width as thermodynamic observable]\label{rem:width-observable}
Resolution width $w(\varphi_n \vdash \bot)$ is robust under time-bounding: unlike Kolmogorov complexity~$K$, the width of a Resolution refutation is a \emph{syntactic} invariant that does not depend on computational resources. This makes it a natural ``thermodynamic observable'' in the proof-complexity analogy: large width forces any time-bounded description to be complex, providing a concrete route from $K$-lower bounds to $K^t$-lower bounds for specific formula families. The conditional MCSP hardness results of Hirahara and Ilango~\cite{Ilango2025} motivate MCSP as a bridge object, without requiring general decidability of~$K$.
\end{remark}

\subsection{\texorpdfstring{Entropy proof of $\SAT \notin \mathrm{AC}^0$}{Entropy proof of SAT not in AC0}}
\label{sec:ac0}

As a concrete instantiation of the entropy framework, we give an entropy-based argument for the classical result $\mathrm{PAR} \notin \mathrm{AC}^0$ (Furst-Saxe-Sipser 1984, Ajtai 1983). Since $\mathrm{AC}^0$ is closed under $\mathrm{AC}^0$-reductions and parity is $\mathrm{AC}^0$-reducible to SAT (via a standard reduction that encodes the parity constraint as a CNF formula using constant-depth circuitry), this also implies $\SAT \notin \mathrm{AC}^0$.

\begin{proposition}[Entropy obstruction for $\mathrm{AC}^0$]
\label{prop:ac0}
Let $\{C_n\}$ be a constant-depth circuit family of polynomial size $s(n) = n^{O(1)}$ computing a Boolean function $f_n: \{0,1\}^n \to \{0,1\}$. Then the truth table $f_n \in \{0,1\}^{2^n}$ satisfies
\begin{equation}
K^{s(n) \cdot 2^n}(f_n) \leq O(n^c \log n)
\end{equation}
for some constant $c$ depending on the depth and size bound. In contrast, a \textit{random} Boolean function $g_n$ satisfies $K(g_n) \geq 2^n - O(1)$ with high probability. The entropy gap between $\mathrm{AC}^0$-computable functions and random functions is thus exponential: $O(n^c \log n)$ vs.\ $2^n$.
\end{proposition}

\begin{proof}[Proof sketch]
A depth-$d$ circuit of size $s$ on $n$ inputs is described by $O(s \cdot \log s)$ bits (gate types and wiring). Given this description, the truth table on all $2^n$ inputs is computed in time $O(s \cdot 2^n)$ by brute-force evaluation. Thus $K^{O(s \cdot 2^n)}(f_n) \leq O(s \cdot \log s) + O(\log n)$. For polynomial size $s = n^c$: $K^{n^c \cdot 2^n}(f_n) \leq O(n^c \log n)$.

The classical result $\mathrm{PAR}_n \notin \mathrm{AC}^0$ (Furst-Saxe-Sipser 1984, Ajtai 1983, H\r{a}stad 1987) shows that the parity function cannot be computed by polynomial-size constant-depth circuits. In the entropy framework, this means: any circuit computing $\mathrm{PAR}_n$ requires super-polynomial size $s(n) = 2^{n^{\Omega(1/d)}}$, so $K^{\mathrm{poly}(2^n)}(\mathrm{PAR}_n)$ is \textit{not} bounded by $O(n^c \log n)$ when computed via $\mathrm{AC}^0$ circuits. The entropy method does not \textit{reprove} $\mathrm{PAR} \notin \mathrm{AC}^0$, but it \textit{reframes} the separation: $\mathrm{AC}^0$-computable truth tables occupy a negligible fraction of the space of all Boolean functions, and parity lies outside this low-entropy region.
\end{proof}

This demonstrates the entropy perspective on a known separation. For general $\PP/\mathrm{poly}$, the analogous bound would give $K^{\mathrm{poly}(2^n)}(f_n) \leq O(n^c \log n)$ for any $f \in \PP/\mathrm{poly}$. Proving that $\SAT_n$ exceeds this bound would establish $\NP \not\subseteq \PP/\mathrm{poly}$ -- a major open problem that implies $\PP \neq \NP$.

\begin{remark}[H\r{a}stad's bound as entropy capacity]
\label{rem:hastad-entropy}
H\r{a}stad's switching lemma~\cite{Hastad1987} can be interpreted as an \emph{information-theoretic capacity bound}: after applying a random restriction that fixes each variable independently with probability $1-p$, a depth-$d$ $\mathrm{AC}^0$ circuit on the surviving variables simplifies to a decision tree of depth $O(\log n)^{d-1}$. In entropy terms: each layer of an $\mathrm{AC}^0$ circuit can ``aggregate'' at most $O(\log n)$ bits of information from its inputs (via the fan-in bound). After $d$ layers, the total information capacity is $O((\log n)^d)$ bits -- far below the $\Omega(n)$ bits needed to compute parity. This capacity bound $O((\log n)^d) \ll n$ is the entropy-theoretic core of the $\mathrm{PAR} \notin \mathrm{AC}^0$ separation.
\end{remark}

\subsection{Geometric Complexity Theory}

Mulmuley's Geometric Complexity Theory (GCT) \cite{Mulmuley2011} approaches $\PP$ vs $\NP$ through algebraic geometry and representation theory. The GCT program seeks \textit{obstructions} -- representation-theoretic invariants that distinguish easy from hard functions.

In the entropy framework, GCT obstructions are a specific type of \textit{algebraic entropy certificate}: they certify that a function has high ``algebraic complexity'' by exhibiting an invariant that polynomial-size circuits cannot match.

The connection to algorithmic positivity: GCT obstructions are \textit{non-natural} (they are not efficiently computable properties of truth tables) and \textit{non-relativizing} (they depend on algebraic structure). They thus satisfy specifications (S1)--(S2). Whether they satisfy (S3) (non-algebrization) is an active research question.


% ===================================================================
% 10. RELATED WORK
% ===================================================================
\section{Related Work}
\label{sec:related}

The connection between Kolmogorov complexity and computational complexity has been studied extensively since the 1990s. The equivalence between time-bounded Kolmogorov complexity and the P vs NP question is \emph{not} new to this paper; our contribution lies in the specific framing as a No-Go theorem via algorithmic positivity, the systematic barrier analysis, and the identification of concrete bridge strategies. We briefly situate our work within this literature.

\textbf{Instance complexity and the P vs NP equivalence.}
Orponen, Ko, Sch\"oning, and Watanabe~\cite{OrponenKoSchoeningWatanabe1994} introduced instance complexity as a measure for the computational difficulty of individual instances: the $t$-bounded instance complexity $\mathrm{ic}_t(x : A)$ is the length of the shortest program that runs in time $t$, decides $x$ correctly with respect to $A$, and makes no mistakes on other inputs. Fortnow~\cite{Fortnow2000} showed explicitly that $\PP \neq \NP$ is equivalent to the statement that the time-bounded instance complexity of SAT is unbounded. The Entropic Separation Conjecture (ESC) of this paper is a reformulation of this known equivalence in the language of conditional witness entropy rather than instance complexity. Our Theorem~\ref{thm:high_entropy} and Proposition~\ref{prop:equiv} are thus new \emph{formulations} of known results, not new results in themselves.

\textbf{Resource-bounded Kolmogorov complexity and circuit lower bounds.}
Allender et al.~\cite{Allender2006} established deep connections between the Minimum Circuit Size Problem (MCSP), resource-bounded Kolmogorov complexity $K^t$, and derandomization. They showed that $K^t$ can encode significant information about circuit complexity: a truth table $f_n$ computed by circuits of size $s$ satisfies $K^{O(s \cdot 2^n)}(f_n) \leq O(s \log s)$ (the time bound $O(s \cdot 2^n)$ accounts for evaluating the circuit on all $2^n$ inputs), linking slice entropy directly to circuit lower bounds. Allender~\cite{Allender2001} surveyed how Kolmogorov complexity and derandomization interact, arguing that understanding the complexity of ``random strings'' (strings with high $K^t$) is central to resolving P vs NP.

\textbf{Uniform lower bounds via probabilistic $K^t$.}
Santhanam~\cite{Santhanam2023} proposed an algorithmic approach to uniform lower bounds using probabilistic time-bounded Kolmogorov complexity ($\mathrm{pK}^{\mathrm{poly}}$). His approach yields \emph{actual new results} (uniform $\mathrm{AC}^0$ lower bounds) from non-trivial sampling algorithms, without relying on hierarchy theorems. The uniformity bridge of our paper (Section~\ref{sec:uniformity-bridge}) addresses the same conceptual gap that Santhanam tackles with concrete algorithmic tools; his work demonstrates that the gap is partially closable in restricted settings.

\textbf{MCSP and NP-hardness.}
Hirahara~\cite{Hirahara2018} showed that the NP-hardness of MCSP under certain reductions would imply new circuit lower bounds, providing a concrete path from $K^t$-based hardness to P vs NP separations. The chain ``hardness of MCSP $\Rightarrow$ circuit lower bounds $\Rightarrow$ truth-table incompressibility $\Rightarrow$ EH'' motivates our bridge strategy B4.

\textbf{Natural proofs and pseudorandomness.}
The natural proofs barrier~\cite{RazborovRudich1997} is closely connected to the non-computability of $K$: a ``natural'' property must be efficiently computable, but high Kolmogorov complexity is not. Our approach exploits this connection systematically. Impagliazzo and Wigderson~\cite{ImpagliazzoWigderson1997} showed that strong circuit lower bounds imply derandomization ($\mathrm{P} = \mathrm{BPP}$), which in Impagliazzo's ``Five Worlds'' framework~\cite{Impagliazzo1995} corresponds to Heuristica or beyond. The failure of EH would place us in Algorithmica (where P = NP), which is widely believed to be false.

\textbf{One-way functions and $K^t$.}
Liu and Pass~\cite{LiuPass2020} proved the landmark result that one-way functions exist if and only if $K^t$ is average-case hard (for $t(n) \geq (1+\varepsilon)n$). This is the precise connection between resource-bounded Kolmogorov complexity and cryptographic assumptions that underlies our Bridge Target~3 and the MCSP internalisation route (Remark~\ref{rem:mcsp-internalisation}).

\textbf{Five Worlds and ESC.}
The ESC can be situated within Impagliazzo's ``Five Worlds'' framework~\cite{Impagliazzo1995}: ESC holds in Heuristica, Pessiland, Minicrypt, and Cryptomania, and is false only in Algorithmica (where P $=$ NP). The min-entropy variant (Remark~\ref{rem:min-entropy}) distinguishes further: in Pessiland, average-case NP-hardness holds but OWFs do not exist, so the distributional ESC holds but the PRG-based Bridge Target~3 is unavailable. This motivates Bridge Targets~1 and~2 as alternatives that do not require cryptographic assumptions.

\textbf{Distinction from prior work.}
This paper does not claim to establish new connections between Kolmogorov complexity and P vs NP -- such connections have been known since the work of Orponen et al.~\cite{OrponenKoSchoeningWatanabe1994} and Fortnow~\cite{Fortnow2000}, and have been developed further by Allender~\cite{Allender2001,Allender2006}, Hirahara~\cite{Hirahara2018}, Liu--Pass~\cite{LiuPass2020}, and Santhanam~\cite{Santhanam2023}. Our contribution is the specific \emph{presentation}: the framing as a No-Go theorem via algorithmic positivity, the systematic barrier analysis (rigorous for $K$, heuristic for $K^t$), the identification of the uniformity bridge as the precise remaining gap, and the connection to proof complexity via the resolution-width conjecture.

% ===================================================================
% 11. DISCUSSION
% ===================================================================
\section{Discussion}
\label{sec:discussion}

\subsection{What has been achieved}

\begin{enumerate}
\item \textbf{Unified presentation.} The known equivalence between $\PP \neq \NP$ and resource-bounded incompressibility of NP-witnesses~\cite{Fortnow2000,OrponenKoSchoeningWatanabe1994} is presented in the unified language of algorithmic positivity and entropic No-Go theorems.

\item \textbf{Barrier analysis.} Arguments based on \emph{unbounded} Kolmogorov complexity $K$ satisfy all five barrier specifications: non-relativizing, non-natural, non-algebrizing, model-independent, and robust. For the \emph{resource-bounded} variant $K^t$ used in the core results, the barrier status is heuristic: $K^t$ is computable, so it does not automatically escape the natural proofs barrier, and formal non-algebrization remains open (see Remark~\ref{rem:barrier-K-vs-Kt}).

\item \textbf{Witness entropy gap.} Conditional on $\PP \neq \NP$, there exist SAT instances whose valid witnesses are not logarithmically compressible within any fixed polynomial time bound and logarithmic constant. The stronger claim that witnesses have near-maximal conditional Kolmogorov complexity is isolated as Conjecture~\ref{conj:near-maximal-witness-entropy}.

\item \textbf{Algorithmic positivity.} The concept identifies two core structural properties (polynomial-time witness extraction and entropic cheapness of witnesses) that $\PP$-decidable NP languages satisfy and NP-complete problems violate, supplemented by heuristic properties (compositional and monotone stability of witnesses) that provide additional separation evidence.
\end{enumerate}

\subsection{What remains open}

The central open problem -- the \textit{uniformity bridge} -- is the conversion of existential entropy bounds (``there exist hard instances'') to universal algorithm bounds (``no algorithm solves all instances''). At a high level, this gap -- between establishing a non-constructive property and closing the loop with a constructive coupling -- appears in other mathematical contexts as well, though the analogy is at the level of proof strategy, not mathematical content:

\begin{center}
\resizebox{\columnwidth}{!}{%
\begin{tabular}{@{}lll@{}}
\hline
Problem & Non-negative quantity & Missing coupling \\
\hline
Hodge & HR-form $\geq 0$ & $\mathrm{AP} \Rightarrow$ algebraic \\
BSD & $\hat{h} \geq 0$ & Higher Gross-Zagier \\
P vs NP & $K^t(w|x)$ high & Uniformity bridge \\
\hline
\end{tabular}
}
\end{center}

\subsection{Limitations and intellectual honesty}

\textbf{This paper does not prove $\PP \neq \NP$.} We state this unambiguously. The contributions are: (1) a clean reformulation that identifies \textit{what} must be shown (witness incompressibility), (2) a barrier analysis arguing \textit{why} this approach is not blocked by known meta-theorems, and (3) concrete research directions (bridge strategies B1--B4) for closing the remaining gap.

\begin{enumerate}
\item \textbf{Non-computability of $K$.} Kolmogorov complexity is not computable, which means the entropy obstruction cannot be directly ``run'' as an algorithm. This is a feature (it bypasses natural proofs) but also a limitation (it makes the argument non-constructive).

\item \textbf{The unique-witness requirement.} The entropy gap argument is strongest for instances with unique witnesses. For instances with many witnesses, a polynomial-time algorithm might select a low-complexity witness even if high-complexity witnesses exist.

\item \textbf{Equivalence, not implication.} The Entropic Separation Conjecture is logically equivalent to $\PP \neq \NP$ (Remark~\ref{rem:esc-equivalence}). Therefore, this paper's No-Go theorem is a \textit{reformulation}, not an independent proof strategy. The value lies in identifying the information-theoretic structure of the problem.

\item \textbf{Barrier analysis: rigorous for $K$, heuristic for $K^t$.} For \emph{unbounded} $K$, the barrier bypass is well-motivated: $K$ is non-computable (defeating natural proofs), oracle-independent (defeating relativization), and non-algebraic (plausibly defeating algebrization). However, the core results use \emph{resource-bounded} $K^t$, which is computable and whose invariance is only up to polynomial slowdown. For $K^t$, the natural proofs barrier is not automatically bypassed, and formal non-algebrization remains open (Remark~\ref{rem:barrier-K-vs-Kt}). Establishing full barrier immunity for $K^t$-based arguments is an important research direction.
\end{enumerate}

\subsection{The Uniformity Bridge: lemma targets}
\label{sec:uniformity-bridge}

The central open problem of this paper is the \emph{Uniformity Bridge}: transferring the entropy separation from unbounded Kolmogorov complexity $K$ (where barrier immunity holds) to resource-bounded complexity $K^t$ (where computational relevance lies). We identify three concrete lemma targets, any one of which would close the bridge.

\paragraph*{Bridge Target 1: instance compression.}
\emph{Statement:} If no polynomial-time algorithm can compress satisfying assignments of $\varphi_n$ below $|\varphi_n|^{1-\varepsilon}$ bits, then $K^{n^c}(w | \varphi_n) \geq |\varphi_n|^{1-\varepsilon}$, or in the strongest unique-witness form $K^{n^c}(w | \varphi_n) \geq |w| - O(\log n)$, for every valid witness $w$.

\emph{Why it suffices:} Instance compression hardness is a well-studied notion in parameterized complexity~\cite{FortnowSanthanam2011,DellMelkebeek2014}. The connection to $K^t$ lower bounds would follow from a simulation argument: if a short program could compute $w$ from $\varphi_n$ in time $n^c$, it would constitute a compression of $\varphi_n$ via the witness.

\paragraph*{Bridge Target 2: proof complexity.}
\emph{Statement:} For any propositional proof system $P$ and tautology $\tau_n$ encoding the unsatisfiability of $\bar{\varphi}_n$, the proof $\pi$ satisfies $K(\pi | \tau_n) \geq |\pi|^{1-\varepsilon}$.

\emph{Why it suffices:} High proof complexity (superpolynomial proof length) implies that proofs themselves carry high entropy. Combined with the Entropic Separation Conjecture, this yields: no polynomial-time proof search can uniformly find proofs, because the proof objects are incompressible. This route connects to the Razborov--Rudich natural proofs barrier in a constructive way.

\paragraph*{Bridge Target 3: PRG-based families.}
\emph{Statement:} There exists an explicit family $\{\varphi_n\}$ of SAT instances such that (i) each $\varphi_n$ has a unique satisfying assignment $w_n$, (ii) $w_n$ is the output of a pseudorandom generator $G: \{0,1\}^{n/2} \to \{0,1\}^n$, and (iii) $K^{n^c}(w_n | \varphi_n) \geq n/2 - O(\log n)$.

\emph{Why it suffices:} If $G$ is a secure PRG (e.g., based on factoring or lattice assumptions), then (iii) follows from the pseudorandomness: any efficient compression of $w_n$ would break $G$. The unique-witness property (i) ensures the entropy gap applies (Corollary~\ref{cor:gap}). The equivalence of OWF existence and average-case hardness of $K^t$~\cite{LiuPass2020} provides a direct cryptographic interpretation: Bridge Target~3 reduces the Uniformity Bridge to a standard cryptographic assumption.

\begin{remark}[Minimal sufficient bridge]
\label{rem:minimal-bridge}
Any \emph{one} of the three targets suffices to upgrade the Entropic Separation Conjecture from an equivalence (reformulation of P $\neq$ NP) to a conditional proof (P $\neq$ NP under a concrete, falsifiable hypothesis). All three targets are themselves major open problems: Target~3 reduces to the existence of secure PRGs (conditional on standard cryptographic \emph{assumptions}, which are unproven). Target~1 connects to instance compression~\cite{Drucker2015}, where existing results are conditional on the non-collapse of the polynomial hierarchy. Target~2 would require superpolynomial Extended Frege lower bounds, which are widely open. None of these targets should be considered ``accessible'' in the sense of being within current proof technology; they represent \emph{concrete formulations} of the uniformity gap rather than tractable problems.
\end{remark}

\begin{remark}[Instance compression as the primary bridge strategy]
\label{rem:compression-bridge}
We elaborate on Bridge Target~1 as the most promising route, following the ``no sparsification unless PH collapses'' programme of Fortnow--Santhanam~\cite{FortnowSanthanam2011} and Dell--van~Melkebeek~\cite{DellMelkebeek2014}.

The argument proceeds in three steps:

\textbf{Step A: Witness-finder $\Rightarrow$ compressor.}
If a polynomial-time algorithm $A$ produces a valid witness $w = A(x)$ for $x \in \mathrm{SAT}$, then $A$ implicitly defines a \emph{compression}: from $x$, one can produce a shorter representation $y = C(x)$ (via the witness) such that a decompressor $D$ recovers enough structure to decide SAT-membership. Concretely, $|y| \leq |A| + O(\log|x|)$ since $A$ is a fixed program. The key point is that this compression is \emph{non-trivial}: it maps a family of $2^{|\varphi_n|}$ possible instances to $O(1)$ bits of advice (the program $A$), which constitutes an exponential compression.

\textbf{Step B: Compression no-go $\Rightarrow$ uniformity.}
The known instance compression impossibility results~\cite{FortnowSanthanam2011,DellMelkebeek2014,Drucker2015} show: if SAT instances of length $n$ can be compressed to $n^{1-\varepsilon}$ bits, then the polynomial hierarchy collapses (NP $\subseteq$ coNP/poly). Applying this to Step~A: if $A$ exists as a polynomial-time witness-finder for \emph{all} instances, then the resulting compression contradicts PH non-collapse. This is precisely the ``positivity no-go'' pattern of FST: not ``we construct hard instances'' but ``every global simplification has catastrophic structural consequences.''

\textbf{Step C: Translation to $K^t$.}
If compression is impossible (under PH non-collapse), then for infinitely many instances $x$, every polynomial-time witness production must fail in the logarithmic-description regime; an additional quantitative argument would be needed to upgrade this to $K^{n^c}(w | x) \geq |w| - O(\log n)$ for all witnesses $w$. This quantitative upgrade is the near-maximal bridge, not a consequence of the basic reformulation.

The remaining gap is the quantifier order: Step~B gives ``for each $A$, there exist hard instances'' (non-uniform), while the Uniformity Bridge needs ``there exist instances hard for all $A$'' (uniform). Closing this gap requires either (i) a \emph{universal} compression impossibility (not just for each compressor individually), or (ii) a diagonalisation argument that uniformises over all polynomial-time algorithms simultaneously.
\end{remark}

\begin{remark}[The Uniformity Bridge as a mixing obstruction]
\label{rem:mixing-obstruction}
The Uniformity Bridge admits a reformulation in the language of statistical mechanics, connecting directly to the Yang--Mills companion paper. Consider the Gibbs measure over witnesses conditioned on an instance $x$:
\[
\pi_x(w) \;\propto\; \mathbf{1}\{V(x,w) = 1\} \cdot e^{-\beta H_x(w)},
\]
where $H_x$ is any energy function on the witness space (e.g., $H_x \equiv 0$ for the uniform measure over valid witnesses). A polynomial-time search algorithm corresponds to a \emph{local Markov dynamics} (Glauber/heat-bath type) on the witness space that mixes rapidly---i.e., reaches a low-entropy region of $\pi_x$ in $\mathrm{poly}(|x|)$ steps.

The Entropic Separation Conjecture (ESC) can then be restated:
\begin{quote}
\emph{For every NP-complete language $L$ and every local Markov dynamics on witnesses, there exist instance families $\{x_n\}$ for which the mixing time to the witness support is $\geq \exp(\Omega(n))$.}
\end{quote}
This is a \textbf{Dobrushin-type obstruction}: the influence matrix $C = (c_{ij})$ of the witness-space Gibbs measure has spectral radius $\rho(C) \geq 1$ for hard instances, preventing Dobrushin uniqueness and hence rapid mixing. A \emph{superficial} structural parallel exists with the Yang--Mills companion paper~\cite{Geiger2026-YM}: there, $\rho(C) < 1$ holds for $\beta < \beta_0$ (strong coupling), yielding a volume-independent mass gap; here, $\rho(C) \geq 1$ is conjectured for NP-hard instance families. However, this parallel is at the level of the Dobrushin criterion formalism, not at the level of physical or mathematical content -- the two settings involve fundamentally different structures.
\end{remark}

\begin{remark}[Algorithmic Downhill Condition]
\label{rem:algorithmic-downhill}
By analogy with the Downhill Flux Condition (DFC) in the turbulence companion paper~\cite{Geiger2026-TU}, define the \emph{Algorithmic Downhill Condition} (ADC): a search process generates a trajectory $w^{(0)}, w^{(1)}, \ldots, w^{(T)}$ of partial witnesses, and a KL-based functional
\[
F_x(w^{(k)}) \;=\; D_{\mathrm{KL}}\!\left(\delta_{w^{(k)}} \,\|\, \pi_x\right)
\]
measures the ``free-energy distance'' to the witness distribution. ADC asserts that $F_x$ is monotonically non-increasing along any polynomial-time trajectory. For NP-complete instances, ADC must be \emph{violated}: the search process requires ``uphill steps'' (local entropy increase) analogous to backscatter in turbulent cascades. The number of required uphill steps provides a lower bound on the computational complexity of the search---connecting entropic separation to the DFC framework.
\end{remark}

\begin{remark}[Connection to the Overlap Gap Property]
\label{rem:ogp-connection}
The mixing obstruction of Remark~\ref{rem:mixing-obstruction} has a precise counterpart in the combinatorial optimization literature: the \emph{Overlap Gap Property} (OGP), introduced by Gamarnik and Sudan~\cite{GamarnikSudan2017} and developed into a systematic theory of algorithmic barriers~\cite{Gamarnik2021}.

The OGP asserts that for random constraint satisfaction problems near the satisfiability threshold, the solution space exhibits \emph{topological disconnectivity}: the overlap (normalised Hamming distance) between any two near-optimal solutions is either very small or very large, with no intermediate values. This is precisely the geometric content of the Dobrushin obstruction $\rho(C) \geq 1$ from Remark~\ref{rem:mixing-obstruction}: the Gibbs measure over witnesses fragments into exponentially many well-separated clusters (the \emph{shattering transition} of Achlioptas and Coja-Oghlan~\cite{AchlioptasCO2008}), and no local Markov dynamics can mix between them in polynomial time.

The chain of implications is:
\[
\begin{aligned}
\text{Shattering}
&\Rightarrow \text{OGP}\\
&\Rightarrow \text{failure of stable algorithms}\\
&\stackrel{?}{\Rightarrow} \text{high } K^t(w \mid x).
\end{aligned}
\]
The first two arrows are established: Gamarnik, Jagannath, and Wein~\cite{GamarnikJW2020} prove that OGP implies failure of low-degree polynomials, Boolean circuits, and Langevin dynamics; Bresler and Huang~\cite{BreslerHuang2022} establish sharp algorithmic phase transitions for random $k$-SAT using the overlap gap property.

The third arrow---\textbf{OGP $\Rightarrow$ high conditional Kolmogorov complexity $K^t(w|x)$}---is the \emph{central open bridge target} of this paper. A proof would proceed: if the solution space has OGP, then any short program extracting a witness $w$ from $x$ must ``cross'' a free-energy barrier between clusters, requiring either superpolynomial time or superlogarithmic description length. This would convert the geometric obstruction (OGP) into an information-theoretic obstruction (high $K^t$), closing the Uniformity Bridge.
\end{remark}

\subsubsection*{Threshold axiom and diagnostic framework}

The Uniformity Bridge admits a precise formulation as a \emph{threshold axiom}---the minimal structural condition that separates existential hardness (``for each $n$, hard instances exist'') from adversarial hardness (``a single procedure generates hard instances for all $n$'').

\begin{axiom}[Uniformity Bridge -- UB-PvNP]\label{ax:ub-pvnp}
Let $\mathcal{W}_n$ denote the set of candidate high-entropy witnesses for SAT instances of size $n$: pairs $(\varphi_n, w)$ with $K^t(w \mid \varphi_n) \geq n^\varepsilon$. A near-maximal, non-uniform bridge target would assert:
\[
\begin{aligned}
&\forall\, \text{algorithm } A,\quad
  \exists^\infty\, n,\quad
  \exists\, \varphi_n \in \mathrm{SAT}_n:\\
&\qquad \text{all witnesses } w \text{ of } \varphi_n\\
&\qquad \text{satisfy } K^t(w \mid \varphi_n) \geq n^\varepsilon.
\end{aligned}
\]
(Note the quantifier $\exists^\infty n$ (``infinitely often''), not $\forall n$ (``almost everywhere''). Theorem~\ref{thm:high_entropy} proves only the weaker logarithmic-threshold version for each fixed $(q,C)$; the displayed $n^\varepsilon$ threshold is a bridge target.)

The \textbf{Uniformity Bridge} demands a \emph{uniform} construction: a single deterministic procedure $\mathcal{U}$ such that, given $1^n$, $\mathcal{U}(1^n)$ produces an instance $\varphi_n$ with the above property for infinitely many $n$, using resources polynomial in $n$.
\end{axiom}

\begin{remark}[Core intuition]\label{rem:ub-intuition}
Non-uniformity is ``hidden advice'': the hard instances $\varphi_n$ exist but their identity depends on the algorithm $A$ in a non-constructive way. The Uniformity Bridge is the attempt to \emph{eliminate this advice}. The hardness of P\,vs\,NP lives precisely at this threshold: if UB succeeds, P$\,\neq\,$NP follows; if UB fails, the failure must manifest as an explicit \emph{advice requirement}---a quantifiable non-uniformity barrier.
\end{remark}

\begin{remark}[Candidate correspondences around UB-PvNP]\label{prop:ub-equivalences}
The following are candidate formulations and neighbouring targets for the Uniformity Bridge:
\begin{enumerate}
\item[(i)] \textbf{Uniform witness construction:} There exists a polynomial-time $\mathcal{U}$ producing high-entropy instances $\varphi_n = \mathcal{U}(1^n)$ for all $n$.
\item[(ii)] \textbf{Uniform circuit lower bounds:} There exists an explicit Boolean function family $(f_n)_{n \geq 1}$, constructible in $\mathrm{P}$, with $K^t(f_n) \geq n^\varepsilon$ (i.e., no small circuits).
\item[(iii)] \textbf{Uniform pseudorandom generators:} There exists a PRG $G: \{0,1\}^{n^\delta} \to \{0,1\}^n$ computable in $\mathrm{P}$ that fools all polynomial-size circuits.
\item[(iv)] \textbf{Parameter-uniform separation:} The constants in the entropy gap $K^t(w \mid \varphi_n) \geq n^\varepsilon$ can be chosen \emph{independently} of $n$: no $n$-dependent thresholds, case distinctions, or lookup tables.
\end{enumerate}
Each of (i)--(iii), if established in a sufficiently strong and uniform form, would imply P$\,\neq\,$NP. The arrows among them are not used here as proved equivalences; they indicate standard bridge routes through MCSP and hardness-vs-randomness, and each route carries its own hypotheses and quantitative losses.
\end{remark}

\begin{remark}[The $n$-leak scanner: diagnostic for non-uniformity]\label{rem:n-leak}
To locate where uniformity breaks, scan the proof for these four signatures:
\begin{enumerate}
\item \textbf{Advice sites:} Any step of the form ``choose for each $n$\ldots'' without an effective selection rule.
\item \textbf{Lookup tables:} Implicit use of an $n$-specific list or encyclopedia (even if it ``merely exists'').
\item \textbf{Diagonalisation without uniformisation:} An object defined by a pure existence argument (Cantor/K\"onig) that yields no uniform construction.
\item \textbf{Coded instance families:} Hardness shifted into a specially prepared family that already contains $n$-dependent advice.
\end{enumerate}
In the current ESC framework, the $n$-leak sits in Theorem~\ref{thm:slice_nogo}: the hard instance $\varphi_n$ is produced by contradiction (``if all instances were easy, $A$ would solve SAT''), yielding existence without construction. Bridge Target~B3 (PRG-based families) is the most direct attempt to close this leak.
\end{remark}

\begin{remark}[No-go mechanisms as hardness localisation]\label{rem:ub-nogo}
UB-PvNP fails if and only if one of the following obstructions applies:
\begin{enumerate}
\item \textbf{Advice unavoidability:} Every candidate uniformisation implicitly requires $a(n)$ bits of advice, and $a(n) \geq n^\delta$ for some $\delta > 0$. This is the P/poly barrier: uniform hardness requires escaping non-uniform computation.
\item \textbf{Parameter explosion:} Every uniformised construction forces parameters to grow superpolynomially in $n$, making the ``uniform'' procedure effectively non-uniform.
\item \textbf{Instance-family trick:} The proof works only on an $n$-dependently chosen family, so the ``solution'' is really a non-uniformity outsourcing.
\end{enumerate}
Each no-go mechanism, if proven, would constitute a \emph{conditional impossibility theorem}: ``P$\,\neq\,$NP cannot be proven by method $X$ because $X$ requires $a(n)$ bits of advice.'' Such results are barrier theorems in the spirit of Baker--Gill--Solovay and Razborov--Rudich, but targeted at the specific uniformity deficit.
\end{remark}

\begin{remark}[Advice-erosion curve]\label{rem:advice-erosion}
The $n$-leak in Remark~\ref{rem:n-leak} can be turned into a quantitative
diagnostic by allowing an explicit advice budget.  For a samplable
yes-distribution $\mathcal{D}_n$, define an informal residual-entropy profile
\[
E_{\mathcal{D}}(n,b,T)
  := \inf_{A:\,|a_n|\leq b}
     H_{\mathrm{res}}\bigl(W \mid X,A(X,a_n),T\bigr),
\]
where $A$ ranges over $T$-time procedures with at most $b$ bits of
$n$-dependent advice.  The intended interpretation is:
$b=0$ is true uniformity, $b=O(\log n)$ encodes only size or regime
information, while $b\geq n^\delta$ represents hidden non-uniform address
information.

The safe easy direction is only diagnostic: if an NP search problem has a
uniform polynomial-time witness selector, then the corresponding
advice-erosion profile collapses already at $b=0$ and $T=\poly(n)$ for the
selector-induced witness choice.  The open bridge target is the opposite
stress test: find SAT/MCSP/OGP families for which
$E_{\mathcal{D}}(n,\operatorname{polylog}(n),\poly(n))$ remains polynomially large.  This
would not by itself prove P$\,\neq\,$NP, but it would localise precisely how
much non-uniform information is needed before the entropy obstruction erodes.
\end{remark}

\begin{remark}[Search-flow conductance]\label{rem:search-flow-conductance}
A complementary diagnostic measures transport rather than advice.  For a SAT
instance $x$, let $\Omega_x$ be a canonical witness-shadow state space and
$S_x\subseteq\Omega_x$ the valid-witness region.  A restricted search class
$\mathcal{C}$ (for example local, Markovian, DPLL-state, or flip-process
models) induces distributions
\[
\mu_0 \to \mu_1 \to \cdots \to \mu_T
\]
on $\Omega_x$.  Informally, $\Phi_{\mathcal{C}}(x)$ is the smallest flow
capacity across a cut separating the initial mass from $S_x$.

For such restricted models one can aim for a standard bottleneck lemma:
if $\mu_0(S_x)$ is negligible and every separating cut has capacity
$\Phi_{\mathcal{C}}(x)\leq \exp(-n^\varepsilon)$, then no
$\poly(n)$-step process in $\mathcal{C}$ can put constant mass on $S_x$.
This is only a restricted-model obstruction.  The genuinely hard
``flow-shadow'' lemma would have to show that any uniform polynomial-time
witness producer for the chosen SAT family admits such a canonical flow
representation without superpolynomial loss.  That flow-shadow lemma is
another precise form of the Uniformity Bridge.
\end{remark}

\begin{lemma}[UB implies separation]\label{lem:ub-separation}
If UB-PvNP (Axiom~\ref{ax:ub-pvnp}) holds, then:
\begin{enumerate}
\item[(a)] P $\neq$ NP (immediate from the existence of uniform hard instances).
\item[(b)] One-way functions exist (via Impagliazzo--Wigderson~\cite{ImpagliazzoWigderson1997}: uniform circuit lower bounds imply PRGs, which imply OWFs).
\item[(c)] The entropy gap $H_{\mathrm{slice}}(\mathrm{SAT}_n) \geq c \cdot 2^{n^\varepsilon}$ holds uniformly.
\end{enumerate}
Conversely, P $\neq$ NP implies UB-PvNP if and only if the hard instances can be made \emph{explicit}---which is exactly the content of the Uniformity Bridge.
\end{lemma}

\subsection{The meta-pattern}

Across all the problems in this research program, the structure is:

\begin{enumerate}
\item \textbf{Identify the positive form} (height, HR-form, capacity, entropy).
\item \textbf{Show one direction} (algebraic $\Rightarrow$ positive, $\PP \Rightarrow$ low entropy).
\item \textbf{Show the hard direction fails for obstructions} (non-algebraic violates positivity, NP-hard has high entropy).
\item \textbf{Close the loop with a composition/coupling theorem.}
\end{enumerate}

For P vs NP, steps 1--3 are established in this paper. Step 4 -- the uniformity bridge -- remains the central challenge.

\subsection{Open directions}

\begin{remark}[Levin's universal search and uniformity collapse]
\label{rem:levin-uniformity}
Levin's universal search algorithm $U$ runs all programs of length $\leq l$ in dovetailed fashion, allocating time $2^{-|p|} \cdot T$ to program $p$. If P $=$ NP, then $U$ finds witnesses in polynomial time (with a large constant). The near-maximal witness-entropy conjecture provides a complementary perspective: if the relevant witnesses satisfy $K^t(w|\varphi) \geq |w| - O(\log n)$, then $U$'s time allocation is dominated by the \emph{correct} program, whose length is essentially $|w|$. The resulting runtime is $2^{|w|} \cdot \mathrm{poly}(n)$---exponential. This is a heuristic optimality picture for Levin search under the stronger near-maximal hypothesis, not a consequence of Theorem~\ref{thm:high_entropy}.
\end{remark}

\begin{remark}[Valiant--Vazirani reduction and few-witness amplification]
\label{rem:valiant-vazirani}
The Valiant--Vazirani theorem reduces SAT to Unique-SAT with high probability via random hash functions. If the near-maximal witness-entropy conjecture holds for the resulting unique-witness instances, then for a random hash $h$ the unique witness $w_h$ of $\varphi \wedge h(w) = 0^k$ satisfies $K^t(w_h | \varphi, h) \geq |w_h| - O(\log n)$ with high probability (over $h$). This ``few-witness amplification'' would transfer the stronger near-maximal variant to generic SAT instances, at the cost of a randomised reduction. Derandomisation via Nisan--Wigderson generators would be an additional bridge target.
\end{remark}

\begin{remark}[Min-entropy variant]
\label{rem:min-entropy}
The near-maximal $K^t$ formulation can be relaxed to a distributional version using min-entropy: for a samplable distribution $\mathcal{D}$ over SAT instances, define $H_\infty^t(\mathcal{D}) = \min_{(\varphi, w) \in \mathrm{supp}(\mathcal{D})} K^t(w | \varphi)$. If $H_\infty^t(\mathcal{D}) \geq |w| - O(\log n)$ for some samplable $\mathcal{D}$, then average-case hardness of NP follows (Impagliazzo's ``Pessiland''). This relaxation is a distributional bridge target for the stronger near-maximal variant and connects to the rich theory of average-case complexity and one-way functions.
\end{remark}

\begin{remark}[Adversarial witness family]
\label{rem:adversarial-witness}
A stronger variant of the ESC constructs an explicit \emph{adversarial} witness family $\{w_n\}$ with a global consistency check: the witnesses are not independently hard but satisfy a collective constraint $\mathcal{C}(w_1, \ldots, w_N) = 1$ that is easy to verify but hard to satisfy. Any algorithm that computes $w_n$ for all $n \leq N$ must ``know'' the global structure $\mathcal{C}$, requiring description length $\Omega(N)$. This approach avoids the single-instance limitations of Corollary~\ref{cor:gap} and is related to the theory of interactive proofs and PCP.
\end{remark}

\begin{remark}[Compression-to-circuit transfer]
\label{rem:compression-circuit}
For restricted circuit classes $\mathcal{C}$ (e.g., $\mathrm{AC}^0$, monotone circuits), the compression-to-circuit transfer takes a sharper form: if $\varphi_n$ has no $\mathcal{C}$-circuit of size $s$ computing its witness, then $K^t_{\mathcal{C}}(w | \varphi_n) \geq \log \binom{2^n}{s}$ where $K^t_{\mathcal{C}}$ denotes the $\mathcal{C}$-restricted Kolmogorov complexity. For $\mathrm{AC}^0$, Razborov--Smolensky lower bounds give $s \geq 2^{n^{\Omega(1)}}$, yielding superpolynomial $K^t_{\mathrm{AC}^0}$ unconditionally. This provides a ``proof of concept'' for Bridge Target~1 in a restricted model.
\end{remark}

\begin{remark}[MCSP internalisation route]
\label{rem:mcsp-internalisation}
Allender and Das (2017) showed that MCSP (Minimum Circuit Size Problem) is hard for SZK under randomised reductions. Liu and Pass~\cite{LiuPass2020} proved that the average-case hardness of $K^t$ (a close relative of MCSP) is \emph{equivalent} to the existence of one-way functions. Combined with Hirahara's~\cite{Hirahara2018} non-black-box worst-case-to-average-case reductions within NP, this provides a route from MCSP hardness to OWF existence. Our framework extends this chain:
\[
\begin{aligned}
\text{MCSP hard}
&\xRightarrow{\text{Liu--Pass}} \text{OWF exist}\\
&\xRightarrow{\text{Imp.--Levin}} \text{avg-case NP hard}\\
&\Rightarrow \PP \neq \NP .
\end{aligned}
\]
The first implication uses the Liu--Pass equivalence between $K^t$ average-case hardness and OWF existence (the connection to MCSP follows from the Allender et al.\ correspondence between MCSP and $K^t$); the second is a standard reduction; the third is immediate (average-case NP-hardness trivially implies $\PP \neq \NP$, since $\PP = \NP$ would make NP easy even on average). Thus, proving average-case hardness of MCSP would, via Liu--Pass plus standard reductions, yield $\PP \neq \NP$.
\end{remark}

\begin{remark}[Ilango's MCSP NP-hardness]\label{rem:ilango-mcsp}
Hirahara and Ilango \cite{Ilango2025} give evidence for MCSP hardness by proving NP-hardness under deterministic quasi-polynomial non-adaptive reductions from well-studied cryptographic and circuit lower-bound assumptions. This does not settle unconditional NP-hardness of MCSP. For the present framework, the relevance is narrower and safer: their conditional result identifies MCSP as a plausible bridge object for turning time-bounded Kolmogorov complexity into standard complexity-theoretic lower bounds.
\end{remark}

\begin{remark}[Quantum extension]
\label{rem:quantum}
The entropy framework extends naturally to quantum computation. A quantum algorithm with $T$ gates on $n$ qubits can be described by $O(T \log T)$ classical bits (encoding the gate sequence), so $K(\mathrm{desc}(U)) \leq O(T \log T)$ for the classical description of the unitary $U$. If $\mathrm{BQP} = \mathrm{QMA}$, then the polynomial-time quantum algorithm would replace the QMA witness, and the analogous witness entropy gap argument would apply: the resource-bounded quantum Kolmogorov complexity of the witness would need to be simultaneously low (by the BQP algorithm) and high (for instances encoding random strings). This suggests $\mathrm{BQP} \neq \mathrm{QMA}$, though making this rigorous requires a quantum analogue of resource-bounded Kolmogorov complexity, which is an active area of research.
\end{remark}

\begin{remark}[Game-theoretic interaction model]
\label{rem:game-theoretic}
The P vs NP question admits a game-theoretic formulation: an entropy maximiser (``Nature'') selects witness distributions to maximise $K^t(w|\varphi)$, while a circuit minimiser (``Algorithm'') selects circuits to minimise computation time. The minimax value of this zero-sum game represents the optimal trade-off between witness entropy and circuit size. A stronger near-maximal variant would assert that Nature can keep this value positive in a scale-sensitive sense for infinitely many $\varphi$. This formulation connects to the FST programme's use of minimax principles in other domains (Yang--Mills, turbulence), but it is not a theorem of the present paper.
\end{remark}

\begin{remark}[Quantum certificates and MIP*]
\label{rem:quantum-certificates}
The near-maximal witness-entropy variant can be translated into the language of quantum interactive proofs: an entropy gap $K^t(w|\varphi) \geq |w| - O(\log n)$ would imply that no polynomial-size quantum circuit can prepare a quantum state $|\psi_w\rangle$ encoding the witness from $|\varphi\rangle$ alone. In the MIP* framework (Natarajan--Wright~\cite{NatarajanWright2019}, Ji--Natarajan--Vidick--Wright--Yuen~\cite{JiNatarajanVidickWrightYuen2021}), entanglement between provers can be viewed as a mechanism for ``distributing'' the witness entropy across non-communicating parties. Whether entanglement can circumvent such a strengthened entropy barrier is an open analogy, not a claim of this framework.
\end{remark}

% Remark on Kolyvagin-type hierarchy removed: the analogy between Euler systems
% (arithmetic geometry of elliptic curves, Galois cohomology, Selmer groups)
% and SAT-witness spaces is too vague to be productive. The only shared property
% -- "hardness propagates from one instance to related instances" -- is not
% specific enough to constitute a research direction.

\begin{remark}[Self-referential perspective -- informal]
\label{rem:self-referential}
A heuristic perspective on why closing the uniformity bridge is hard: for any fixed polynomial-time machine $A$ of description length $|A|$ and running-time polynomial $q_A$, Theorem~\ref{thm:high_entropy} (conditional on $\PP \neq \NP$) gives, for every sufficiently large logarithmic constant $C$, infinitely many instances $\varphi$ whose valid witnesses all satisfy $K^{q_A} (w \mid \varphi) > C\log|\varphi|$. Since $A$ is a fixed program, $K^{q_A}(A(x) \mid x) \leq |A| + O(\log |x|) \leq C\log |x|$ for all large $x$. These two bounds are incompatible on the corresponding hard instances: $A$ cannot produce a valid witness there. This is simply the non-uniform direction of the witness entropy gap (already established in Theorem~\ref{thm:high_entropy}) and does \emph{not} close the uniformity bridge, which requires hard instances that resist \emph{all} algorithms simultaneously, not just each algorithm individually.
\end{remark}

\begin{remark}[Extended Frege and MCSP]
\label{rem:frege-mcsp}
Two concrete targets for unconditional progress:
\begin{enumerate}
\item \emph{Extended Frege:} If superpolynomial lower bounds on Extended Frege proof length can be established for tautologies encoding the unsatisfiability of ESC-hard instances, then the proof objects themselves have high entropy ($K(\pi | \tau) \geq |\pi|^{1-\varepsilon}$), realising Bridge Target~2.
\item \emph{MCSP $\notin$ P/poly:} The Minimum Circuit Size Problem (Allender--Hirahara~\cite{AllenderHirahara2019}) is the natural decision version of Kolmogorov complexity for circuits. Proving MCSP $\notin$ P/poly would imply uniform entropy hardness (EH) and, via our framework, P $\neq$ NP. This is considered one of the most accessible unconditional targets in complexity theory.
\end{enumerate}
\end{remark}

\begin{remark}[Loose analogy with the Hodge conjecture]
\label{rem:hodge-analogy}
At a very high level, the ESC framework shares a structural pattern with the Hodge conjecture: objects that admit structured descriptions (polynomial-time algorithms / algebraic cycles) form a proper subset of a larger space, and the question is whether certain distinguished objects (NP-witnesses / Hodge classes) lie in this subset. However, this analogy is \emph{superficial}: the Hodge-Riemann bilinear form is a deep topological invariant on cohomology, while $K^t \geq 0$ is a trivial property of any length function. The analogy should be understood as a heuristic parallel in proof strategy, not as evidence for a mathematical connection.
\end{remark}


% ===================================================================
\subsection{Post-Zookeeper transfer: Williams-normalised uniformity}
\label{subsec:zookeeper-pnp-transfer}

The post-Zookeeper role of the present paper is to make the uniformity
gap operational.  The relevant recent input is Williams' square-root-space
simulation of time-bounded computation: time $t$ can be simulated in
space $O(\sqrt{t\log t})$ for the relevant Turing-machine model.  This
does not resolve P versus NP, but it changes the geometry of the
resources that the entropy obstruction must survive.

Relative to the current CORE status, the Zookeeper closure of
Riemann-$C2/MS2$ is a \emph{normal-form constraint}, not a transfer proof
for complexity theory.  It licenses the five-slot bookkeeping below, but
it does not close the uniformity bridge.  The open mathematical step in
this paper remains the passage from non-uniform high-entropy witnesses to
a uniform lower bound that survives the best available time--space
simulation.

The transfer normal form is:
\begin{description}
  \item[Operator.] Uniform search and verification operators, measured by
  resource-bounded Kolmogorov complexity $K^t$.
  \item[Defect.] The uniformity defect: high-entropy witness existence
  minus uniform generation under the best available time--space
  simulation.
  \item[Approximation.] SAT slices, resolution-width regimes, MCSP/XOR/MUX
  candidates, and circuit-size levels.
  \item[Gap.] A residual entropy hardness gap after time has been
  normalised by the square-root-space simulation.
  \item[Limit.] Passage from non-uniform instance families to a uniform
  language or algorithm.
\end{description}

This suggests a refined diagnostic:
\[
\begin{aligned}
  H_{\mathrm{slice}}^{\mathrm{ts}}(n;s)
  &= \text{witness entropy after simulation}\\
  &\quad \text{in space }s,\qquad s \asymp \sqrt{t\log t}.
\end{aligned}
\]
The next target is not a direct proof of $\mathrm{P}\ne\mathrm{NP}$, but a
transfer lemma: if a uniform polynomial-time search procedure exists, then
this Williams-normalised slice entropy must collapse along the associated
resource surface.  MCSP-style candidates and resolution-width families are
then the natural stress tests for non-collapse.

\subsection*{Result classification}

This paper is a \textbf{framework and presentation paper}: it presents the known equivalence between the P\,$\neq$\,NP conjecture and an entropic condition (ESC) in a unified framework, without proving the separation itself.

\begin{table*}[t]
\caption{Result classification.}
\label{tab:result-classification}
\centering
\footnotesize
\resizebox{\textwidth}{!}{%
\begin{tabular}{@{}lll@{}}
\toprule
\textbf{Result} & \textbf{Type} & \textbf{Reference} \\
\midrule
\multicolumn{3}{@{}l}{\textit{Proven equivalences:}} \\
ESC $\Leftrightarrow$ P\,$\neq$\,NP & Proposition + Remark & Prop.~\ref{prop:equiv}, Rem.~\ref{rem:esc-equivalence} \\
EH $\Rightarrow$ P\,$\neq$\,NP & Theorem & Thm.~\ref{thm:slice_nogo} \\
\midrule
\multicolumn{3}{@{}l}{\textit{Structural contributions:}} \\
Barrier immunity (S1--S5 for $K$; heuristic for $K^t$) & Analysis & Sec.~\ref{sec:barriers}, Rem.~\ref{rem:barrier-K-vs-Kt} \\
Resolution-width conjecture & Conjecture & Conj.~\ref{conj:resolution-width} \\
MCSP internalisation route & Remark & Rem.~\ref{rem:mcsp-internalisation} \\
Advice-erosion curve & Bridge target & Rem.~\ref{rem:advice-erosion} \\
Search-flow conductance & Bridge target & Rem.~\ref{rem:search-flow-conductance} \\
\midrule
\multicolumn{3}{@{}l}{\textit{Exploratory numerics (not constituting evidence):}} \\
Phase transition structure observed & Illustrative & Rem.~\ref{rem:sat-entropy} \\
$H_{\mathrm{slice}}$ growth trend ($n \leq 16$, 5 points) & Illustrative & Rem.~\ref{rem:sat-entropy} \\
\midrule
\multicolumn{3}{@{}l}{\textit{Open (equivalent to P\,$\neq$\,NP):}} \\
Uniformity Bridge & Open problem & Sec.~\ref{sec:uniformity-bridge} \\
Entropy Hardness (EH) & Conjecture & Sec.~\ref{sec:slices} \\
\bottomrule
\end{tabular}
}
\end{table*}

% ===================================================================
% 12. CONCLUSION
% ===================================================================
\section{Conclusion}
\label{sec:conclusion}

The equivalence between $\PP \neq \NP$ and the incompressibility of NP-witnesses (under resource-bounded Kolmogorov complexity) has been known in various forms since the work of Orponen et al.~\cite{OrponenKoSchoeningWatanabe1994} and Fortnow~\cite{Fortnow2000}. This paper presents these known connections in a unified framework, introducing the terminology of algorithmic positivity and the entropic No-Go theorem.

The witness entropy gap -- the discrepancy between the high conditional complexity of witnesses (under $\PP \neq \NP$) and the low complexity that polynomial-time algorithms can produce -- provides a clean characterization of what any proof of $\PP \neq \NP$ must establish. The barrier analysis argues that $K$-based arguments are not blocked by the three classical barriers, though this immunity is rigorous only for unbounded $K$ and remains heuristic for the resource-bounded $K^t$ used in the core results.

The remaining challenge is the uniformity bridge: converting the existential entropy bound into a universal lower bound against all polynomial-time algorithms. The bridge strategies identified (instance compression, proof complexity, PRG-based families, MCSP) are themselves major open problems, each essentially equivalent to or harder than P vs NP itself. The value of the framework lies in making this structure explicit.


% ===================================================================
% ACKNOWLEDGMENTS
% ===================================================================
\begin{acknowledgments}
This work was conducted as independent research without institutional funding.

\textit{AI tools.} AI-based tools (Claude by Anthropic) were used for mathematical formalization and text generation. All content and responsibility for correctness lie solely with the author.

\textit{Funding information.} This research received no external funding.
\end{acknowledgments}


% ===================================================================
% BIBLIOGRAPHY
% ===================================================================
\begin{thebibliography}{33}

\bibitem{BakerGillSolovay1975}
T.~Baker, J.~Gill, \& R.~Solovay (1975).
Relativizations of the $\mathrm{P} =\mathrel{?} \mathrm{NP}$ question.
\textit{SIAM Journal on Computing}, 4(4), 431--442.
DOI: 10.1137/0204037.

\bibitem{RazborovRudich1997}
A.~Razborov \& S.~Rudich (1997).
Natural proofs.
\textit{Journal of Computer and System Sciences}, 55(1), 24--35.
DOI: 10.1006/jcss.1997.1494.

\bibitem{AaronsonWigderson2009}
S.~Aaronson \& A.~Wigderson (2009).
Algebrization: A new barrier in complexity theory.
\textit{ACM Transactions on Computation Theory}, 1(1), 2.
DOI: 10.1145/1490270.1490272.

\bibitem{LiVitanyi2008}
M.~Li \& P.~Vit\'anyi (2008).
\textit{An Introduction to Kolmogorov Complexity and Its Applications}, 3rd ed.
Springer.
DOI: 10.1007/978-0-387-49820-1.

\bibitem{Cook1975}
S.~A.~Cook (1975).
Feasibly constructive proofs and the propositional calculus.
\textit{Proceedings of the 7th ACM Symposium on Theory of Computing}, pp.\ 83--97.
DOI: 10.1145/800116.803756.

\bibitem{BenSasson2001}
E.~Ben-Sasson \& A.~Wigderson (2001).
Short proofs are narrow -- resolution made simple.
\textit{Journal of the ACM}, 48(2), 149--169.
DOI: 10.1145/375827.375835.

\bibitem{Razborov1985}
A.~Razborov (1985).
Lower bounds on the monotone complexity of some Boolean functions.
\textit{Doklady Akademii Nauk SSSR}, 281(4), 798--801.

\bibitem{Razborov1987}
A.~Razborov (1987).
Lower bounds on the size of bounded depth circuits over a complete basis with logical addition.
\textit{Mathematical Notes}, 41(4), 333--338.
DOI: 10.1007/BF01137685.

\bibitem{Grigoriev2001}
D.~Grigoriev (2001).
Linear lower bound on degrees of Positivstellensatz calculus proofs for the parity.
\textit{Theoretical Computer Science}, 259(1--2), 613--622.
DOI: 10.1016/S0304-3975(00)00403-7.

\bibitem{Mulmuley2011}
K.~D.~Mulmuley (2011).
On P vs.\ NP and geometric complexity theory.
\textit{Journal of the ACM}, 58(2), 5.
DOI: 10.1145/1667053.1667055.

\bibitem{Arora2009}
S.~Arora \& B.~Barak (2009).
\textit{Computational Complexity: A Modern Approach}.
Cambridge University Press.

\bibitem{Kolmogorov1965}
A.~N.~Kolmogorov (1965).
Three approaches to the quantitative definition of information.
\textit{Problems of Information Transmission}, 1(1), 1--7.

\bibitem{Chaitin1966}
G.~J.~Chaitin (1966).
On the length of programs for computing finite binary sequences.
\textit{Journal of the ACM}, 13(4), 547--569.
DOI: 10.1145/321356.321363.

\bibitem{Monasson1999}
R.~Monasson, R.~Zecchina, S.~Kirkpatrick, B.~Selman, \& L.~Troyansky (1999).
Determining computational complexity from characteristic `phase transitions'.
\textit{Nature}, 400, 133--137.
DOI: 10.1038/22055.

\bibitem{ImpagliazzoWigderson1997}
R.~Impagliazzo \& A.~Wigderson (1997).
P = BPP if E requires exponential circuits: Derandomizing the XOR Lemma.
\textit{Proceedings of the 29th ACM STOC}, pp.\ 220--229.
DOI: 10.1145/258533.258590.

\bibitem{FST-Unified2026}
L.~Geiger (2026).
The Renormalized Free Energy Framework.
Zenodo preprint.
\href{https://doi.org/10.5281/zenodo.19036190}{doi:10.5281/zenodo.19036190}.

\bibitem{GeigerRH2026c}
L.~Geiger (2026).
The Riemann Hypothesis Part III: The Analytical Rank-Finite Certificate.
Zenodo preprint.
\href{https://doi.org/10.5281/zenodo.19035640}{doi:10.5281/zenodo.19035640}.

\bibitem{Allender2006}
E.~Allender, H.~Buhrman, M.~Kouck\'y, D.~van~Melkebeek, \& D.~Ronneburger (2006).
Power from Random Strings.
\textit{SIAM Journal on Computing}, 35(6), 1467--1493.
DOI: 10.1137/050628994.

\bibitem{Hastad1987}
J.~H\r{a}stad (1987).
\textit{Computational Limitations of Small-Depth Circuits}.
MIT Press.

\bibitem{Cook1971}
S.~A.~Cook (1971).
The complexity of theorem-proving procedures.
\textit{Proceedings of the 3rd ACM Symposium on Theory of Computing}, pp.\ 151--158.
DOI: 10.1145/800157.805047.

\bibitem{ValiantVazirani1986}
L.~Valiant \& V.~Vazirani (1986).
NP is as easy as detecting unique solutions.
\textit{Theoretical Computer Science}, 47, 85--93.
DOI: 10.1016/0304-3975(86)90135-0.

\bibitem{Solomonoff1964}
R.~J.~Solomonoff (1964).
A formal theory of inductive inference, Parts I and II.
\textit{Information and Control}, 7(1), 1--22 and 7(2), 224--254.
DOI: 10.1016/S0019-9958(64)90131-7.

\bibitem{Allender2001}
E.~Allender (2001).
When worlds collide: Derandomization, lower bounds, and Kolmogorov complexity.
\textit{Proceedings of the 21st FSTTCS}, LNCS 2245, pp.\ 1--15.
DOI: 10.1007/3-540-45294-X\_1.

\bibitem{Fortnow2000}
L.~Fortnow (2000).
Kolmogorov complexity and computational complexity.
Lecture notes, Mathematics Workshop, Kaikoura, New Zealand; published as survey in \textit{Quaderni di Matematica}.

\bibitem{OrponenKoSchoeningWatanabe1994}
P.~Orponen, K.-I.~Ko, U.~Sch\"oning, \& O.~Watanabe (1994).
Instance complexity.
\textit{Journal of the ACM}, 41(1), 96--121.
DOI: 10.1145/174644.174648.

\bibitem{Santhanam2023}
R.~Santhanam (2023).
An algorithmic approach to uniform lower bounds.
\textit{Proceedings of the 38th Computational Complexity Conference (CCC 2023)}, LIPIcs Vol.\ 264, Article~35.
DOI: 10.4230/LIPIcs.CCC.2023.35.

\bibitem{Hirahara2018}
S.~Hirahara (2018).
Non-black-box worst-case to average-case reductions within NP.
\textit{Proceedings of the 59th IEEE FOCS}, pp.\ 247--258.
DOI: 10.1109/FOCS.2018.00032.

\bibitem{Impagliazzo1995}
R.~Impagliazzo (1995).
A personal view of average-case complexity.
\textit{Proceedings of the 10th Structure in Complexity Theory Conference}, pp.\ 134--147.
DOI: 10.1109/SCT.1995.514853.

\bibitem{AchlioptasRicci2006}
D.~Achlioptas, F.~Ricci-Tersenghi (2006).
On the solution-space geometry of random constraint satisfaction problems.
\textit{Proceedings of the 38th Annual ACM Symposium on Theory of Computing (STOC '06)}, pp.\ 130--139.
DOI: 10.1145/1132516.1132537.

\bibitem{MezardZecchina2002}
M.~M\'ezard, R.~Zecchina (2002).
Random $K$-satisfiability problem: From an analytic solution to an efficient algorithm.
\textit{Physical Review E}, 66, 056126.
DOI: 10.1103/PhysRevE.66.056126.

\bibitem{FortnowSanthanam2011}
L.~Fortnow \& R.~Santhanam (2011).
Infeasibility of instance compression and succinct PCPs for NP.
\textit{Journal of Computer and System Sciences}, 77(1), 91--106.
DOI: 10.1016/j.jcss.2010.06.007.

\bibitem{DellMelkebeek2014}
H.~Dell \& D.~van~Melkebeek (2014).
Satisfiability allows no nontrivial sparsification unless the polynomial-time hierarchy collapses.
\textit{Journal of the ACM}, 61(4), Article~23.
DOI: 10.1145/2629620.

\bibitem{Drucker2015}
A.~Drucker (2015).
New limits to classical and quantum instance compression.
\textit{SIAM Journal on Computing}, 44(5), 1443--1479.
DOI: 10.1137/130927115.

\bibitem{AllenderHirahara2019}
E.~Allender \& S.~Hirahara (2019).
New insights on the (non-)hardness of circuit minimization and related problems.
\textit{ACM Transactions on Computation Theory}, 11(4), 27:1--27:27.
DOI: 10.1145/3349616.

\bibitem{NatarajanWright2019}
A.~Natarajan \& J.~Wright (2019).
NEEXP $\subseteq$ MIP*.
\textit{Proceedings of the 60th IEEE FOCS}, pp.\ 510--528.
DOI: 10.1109/FOCS.2019.00039.

\bibitem{JiNatarajanVidickWrightYuen2021}
Z.~Ji, A.~Natarajan, T.~Vidick, J.~Wright, \& H.~Yuen (2021).
MIP* = RE.
\textit{Communications of the ACM}, 64(11), 131--138.
DOI: 10.1145/3485628.

\bibitem{Ilango2025}
S.~Hirahara and R.~Ilango,
\textit{NP-hardness of the Minimum Circuit Size Problem from Well-Studied Assumptions},
FOCS Proceedings, 2025.

\bibitem{GamarnikSudan2017}
D.~Gamarnik and M.~Sudan, \emph{Limits of local algorithms over sparse random graphs}, Ann.\ Probab.\ \textbf{45} (2017), 2353--2376.

\bibitem{Gamarnik2021}
D.~Gamarnik, \emph{The overlap gap property: a topological barrier to optimizing over random structures}, Proc.\ Natl.\ Acad.\ Sci.\ \textbf{118}(41) (2021), e2108492118.

\bibitem{AchlioptasCO2008}
D.~Achlioptas and A.~Coja-Oghlan, \emph{Algorithmic barriers from phase transitions}, Proc.\ 49th FOCS (2008), 793--802.

\bibitem{GamarnikJW2020}
D.~Gamarnik, A.~Jagannath, and A.~S.~Wein, \emph{Hardness of random optimization problems for Boolean circuits, low-degree polynomials, and Langevin dynamics}, preprint, arXiv:2004.12063 (2020).

\bibitem{BreslerHuang2022}
G.~Bresler and B.~Huang, \emph{The algorithmic phase transition of random $k$-SAT for low degree polynomials},
\textit{Proceedings of the 63rd IEEE FOCS} (2022), pp.\ 298--309.
arXiv:2106.02129.

\bibitem{LiuPass2020}
Y.~Liu and R.~Pass (2020).
On one-way functions and Kolmogorov complexity.
\textit{Proceedings of the 61st IEEE FOCS}, pp.\ 1243--1254.
DOI: 10.1109/FOCS46700.2020.00119. arXiv:2009.11514.

\bibitem[Geiger(2026d)]{Geiger2026-TU}
Geiger, L. (2026).
\newblock Functional Stability Theory: Turbulence and the K41 Spectrum as Free-Energy Minimum.
\newblock Preprint (FST-TU).

\bibitem[Geiger(2026e)]{Geiger2026-YM}
Geiger, L. (2026).
\newblock Functional Stability Theory: Volume-Independent Spectral Gap for Lattice Yang--Mills.
\newblock Preprint (FST-YM).

\end{thebibliography}

\end{document}
